\documentclass[11pt,a4paper]{article}
\usepackage[T1]{fontenc}
\usepackage{lmodern}
\usepackage[utf8]{inputenc}
\usepackage{amsmath}
\usepackage{amssymb}
\usepackage{amsthm}
\usepackage{booktabs}
\usepackage{geometry}
\usepackage{xcolor}
\usepackage{enumitem}
\usepackage{array}
\usepackage{graphicx}
\usepackage{microtype}
\usepackage{hyperref}
\usepackage{xurl}

\geometry{margin=0.85in}
\hypersetup{colorlinks=true,linkcolor=blue!70!black,urlcolor=blue!70!black,citecolor=blue!70!black}

\newcommand{\spath}[1]{\path{#1}}
\newcommand{\fref}[1]{[\path{#1}]}
\setlength{\emergencystretch}{4em}
\sloppy
\hbadness=10000
\hfuzz=2pt
\numberwithin{equation}{section}
\numberwithin{table}{section}

\newtheorem{theorem}{Theorem}[section]
\newtheorem{lemma}[theorem]{Lemma}
\newtheorem{proposition}[theorem]{Proposition}
\newtheorem{corollary}[theorem]{Corollary}
\newtheorem{definition}[theorem]{Definition}
\newtheorem{conjecture}[theorem]{Conjecture}
\newtheorem{remark}[theorem]{Remark}
\newtheorem{assumption}[theorem]{Assumption}

\title{\textbf{Diamond Framework:}\\[0.5em]Reputation-Weighted Validator Selection\\via Verifiable Non-Transferable Reputation}
\author{Ilich Blanco\thanks{Andromeda Computer, \texttt{ilichblanco@andromedacomputer.net}}}
\date{September 2026}

\begin{document}
\maketitle

\begin{abstract}
The Diamond Framework addresses the structural tension in Proof-of-Stake between security and accessibility. High collateral requirements deter honest but capital-constrained validators; low requirements expose the network to Sybil attacks. The framework replaces a fraction of the financial stake with verifiable, non-transferable reputation.

The paper makes four contributions: a cryptoeconomic protocol in which financial stake and reputational collateral are partial substitutes; a numerical verification suite that pairs every quantitative claim with a reproducible script; a falsification log documenting twenty-one claims that were made, tested, and corrected during development; and a set of design decisions with explicit boundaries, stating where the protocol cannot deliver a guarantee.

The protocol is verified numerically for the reputation dynamics, the tensor representation, the zero-knowledge threshold proof, the collusion detection, the payment rule, the data layer, and the bootstrap transition. EVM precompile costs are measured locally with Foundry. Two benchmarks are deferred: the recursive SNARK proving time and the Teleporter latency, both requiring infrastructure not available in this paper.
\end{abstract}

\tableofcontents
\newpage

\section{Introduction}

\subsection{The Validator Trust Gap}
Proof-of-stake networks secure their consensus by requiring validators to lock financial collateral. Slashing conditions make misbehavior costly, and the resulting cryptoeconomic security guarantees are well formalized. This mechanism has proven effective. It has also produced a structural tension.

High stake requirements deter honest but capital-constrained participants, concentrating validation power in the hands of large token holders. Low requirements expose the network to Sybil attacks and near-zero-cost misbehavior. The tension is not a design flaw of any particular protocol; it is a consequence of using a single instrument---financial capital---to serve two purposes simultaneously: deterrence and admission.

In proof-of-work networks, physical capacity constraints (hardware, energy, cooling) act as a natural brake on centralization. A miner with a technological advantage cannot expand indefinitely without facing diminishing returns on capacity. In proof-of-stake networks the collateral is purely financial and does not face an equivalent physical brake. The resulting concentration dynamics can be self-reinforcing.

\subsection{Reputation as Partial Collateral}
The Diamond Framework introduces a second instrument to the admission decision: verifiable, non-transferable reputation. A validator with a long history of contributions faces a large opportunity cost if it misbehaves, because destruction of its reputation eliminates future income streams. This opportunity cost can serve as partial collateral, reducing the financial barrier for honest participants without degrading the Cost of Attack that secures the network.

Making this intuition rigorous requires solving several problems:

\begin{enumerate}[leftmargin=*]
\item \textbf{Stochastic modeling.} Reputation evolves continuously and is subject to abrupt downward shocks. The dynamics must remain in their positive domain under all admissible shocks.
\item \textbf{Verifiability.} Reputation claims must be cryptographically verifiable on-chain and without relying on a centralized oracle. The tensor $T$ is public (Section~\ref{sec:datalayer}), so the proof provides succinctness rather than privacy.
\item \textbf{Collusion resistance.} Validators may coordinate to inflate each other's reputation. The detection mechanism must be robust to heavy-tailed interaction graphs and must have a stated boundary below which detection is information-theoretically impossible.
\item \textbf{Incentive compatibility.} The mechanism must ensure validators report their true reputation, or---if reputation is publicly verified---must make the reporting question vacuous.
\item \textbf{Data integrity.} The reputation signal is derived from consensus events on source chains. The transport of these events must be verifiable, and fraud must be detectable and deterred.
\end{enumerate}

The Diamond Framework addresses all five.

\subsection{Contributions}
The paper makes four contributions.

\textbf{A cryptoeconomic protocol.} We present a validator admission rule in which financial stake and reputational collateral are partial substitutes. The rule is $\mathrm{CoA}(i) = S_f(i) + C_R(T(i)) \geq \Theta(t)$, where $C_R(T)$ is the reputational collateral derived from a verified reputation tensor $T$. The protocol includes the reputation dynamics, the collusion test, the payment rule, the data layer, and the bootstrap simulation.

\textbf{A numerical verification suite.} Every quantitative claim in the paper is paired with a script that reproduces it. The scripts are self-contained, depend only on NumPy and SciPy, and declare their seeds and sample sizes. The suite covers stochastic simulation, spectral calibration, game-theoretic equilibrium, and EVM gas benchmarks. Claims without scripts are labeled as conjectures.

\textbf{A falsification log.} During development, twenty-one quantitative claims were made, tested, and either confirmed or corrected. The corrections include a fabricated simulation figure, a mismeasurement of a gas cost, a factor of two in an integral, a rounding rule that shifted a phase transition, and an unverified projection of a multidimensional type. The log is part of the paper's record, not a supplement to it.

\textbf{Design decisions with explicit boundaries.} Where the protocol cannot deliver a guarantee---for adversaries that are not rational, for collusion below the Kesten--Stigum boundary, for real reputation data that is not available---the paper states the boundary and does not claim the guarantee. The limitations section separates methodological limits (the numerical methods could be improved) from design limits (the protocol does not solve these problems even with perfect methods).

\subsection{Paper Structure}
Section~\ref{sec:reputation} presents the reputation dynamics: a CIR diffusion with multiplicative downward jumps, the positivity proof, and the exact mean first-passage time.

Section~\ref{sec:tensor} introduces the tensor representation and the Tucker compression, and shows that the compression fails for block structure while remaining useful for low-rank tensors.

Section~\ref{sec:zk} specifies the zero-knowledge threshold proof. The original circuit is unsound; the repair uses a Pedersen commitment with Bulletproofs and recursive SNARK composition. The EVM precompile costs are measured locally.

Section~\ref{sec:collusion} develops the collusion detection. The modularity matrix $B = W - dd^{\top}/(2m)$ replaces the normalized adjacency matrix $A_{\mathrm{sym}}$. The null quantile is calibrated by Monte Carlo, and the detection boundary is characterized.

Section~\ref{sec:payment} specifies the payment rule: a saturated uniform posted price. Incentive Compatibility is tautological because reputation is verified on-chain and the allocation does not depend on any private type.

Section~\ref{sec:datalayer} describes the data layer: the ingestion of consensus events, the transport via state inclusion proofs, and the fraud proof mechanism with its audit equilibrium.

Section~\ref{sec:bootstrap} presents the bootstrap simulation under the net\_gain reputation model, with the comparison of three policies for the transition from pure proof-of-stake to hybrid.

Section~\ref{sec:evaluation} evaluates reproducibility, on-chain gas costs, and the stability of the conclusions under parameter variation.

Section~\ref{sec:limitations} states the limitations of the design and of the numerical verification, separated by category.

Section~\ref{sec:conclusion} concludes.

\subsection{Reproducibility}
The paper is accompanied by a public repository containing:

\begin{itemize}[leftmargin=*]
\item The verification scripts (\spath{verification/}).
\item The EVM benchmarks (\spath{integrations/avalanche/benchmarks/}).
\item The falsification log (\spath{docs/audit/numerical-verification.md}).
\item The paper source (\spath{docs/paper/main.tex}).
\end{itemize}

Every quantitative claim in the paper can be traced to a script, a benchmark, or a log entry. The commands to reproduce the claims are listed in \spath{verification/README.md}. Running the full suite takes approximately 67 minutes on a single core, dominated by the heavy-tail recalibrations (\spath{heavy_tail_n5000.py}, $\sim$15 min; \spath{heavy_tail_scaled_clique.py}, $\sim$24 min) and the RMT calibration (\spath{recalibrate_top.py}, $\sim$10 min).

\section{Reputation Dynamics}
\label{sec:reputation}

\subsection{The CIR Process}
Validator reputation is modeled as a mean-reverting stochastic process with abrupt downward shocks. In its continuous form, the process follows the Cox--Ingersoll--Ross (CIR) diffusion:

\begin{equation}
dR(t) = \alpha(\bar{R} - R(t))\,dt + \sigma_0\sqrt{R(t)}\,dW(t)
\label{eq:cir}
\end{equation}

with $\alpha > 0$ the mean-reversion speed, $\bar{R} > 0$ the long-run mean, and $\sigma_0 > 0$ the diffusion coefficient. The process is strictly positive under the Feller condition $2\alpha\bar{R} \geq \sigma_0^2$, which we assume throughout.

Slashing events are modeled as multiplicative downward jumps:

\begin{equation}
R(t) \mapsto R(t^-) \cdot y, \quad y \in (0, 1]
\label{eq:jumps}
\end{equation}

where $y$ is drawn from a L\'evy measure $\nu$ supported on $(0, 1]$. The choice of multiplicative rather than additive jumps is not cosmetic; it is required for the process to remain in its domain.

\subsection{Domain Validity Under Jumps}
The original formulation of the Diamond Framework used additive jumps of the form:

\begin{equation}
dR(t) = \alpha(\bar{R} - R(t))\,dt + \sigma_0\sqrt{R(t)}\,dW(t) + \int J(z)\tilde{N}(dt, dz)
\label{eq:additive}
\end{equation}

with $J(z) < 0$. This specification is not well posed. Additive negative jumps can push $R$ below zero regardless of the Feller condition, because the Feller condition constrains the diffusion, not the jump component.

We quantify this failure by Monte Carlo simulation. With parameters $\alpha = 0.5$, $\bar{R} = 1.0$, $\sigma_0 = 0.8$ (Feller satisfied: $2\alpha\bar{R} = 1.0 \geq \sigma_0^2 = 0.64$), jump intensity $\lambda = 0.8$, mean jump magnitude $0.15$, $R(0) = 0.3$, and $T = 5$, the additive-jump specification produces:

\begin{center}
\begin{tabular}{lr}
\toprule
Specification & Fraction with $R(t) < 0$ \\
\midrule
Additive jumps & 40.58\% \\
Multiplicative jumps & 0.0000\% \\
\bottomrule
\end{tabular}
\end{center}

over 100{,}000 simulated paths with $dt = 0.001$ \fref{verification/stochastics/sim_positivity.py}.

The failure rate of 40.58\% is not a tail event. With the specified parameters, two in five trajectories violate the domain within five time units. The multiplicative specification resolves this by construction: a jump of the form $R \mapsto R \cdot y$ with $y \in (0, 1]$ preserves positivity pathwise, and the CIR component remains positive under Feller.

\subsection{The Whitewashing Bound}
For a validator whose reputation has been reduced to near zero, the expected time to recover to a target level $R_0$ is governed by the first-passage time of the CIR diffusion. Under the scale function $s'(x) = x^{-\kappa}\exp(\beta x)$ with $\kappa = 2\alpha\bar{R}/\sigma_0^2$ and $\beta = 2\alpha/\sigma_0^2$, the classical expression for the expected first-passage time is:

\begin{equation}
\mathbb{E}[\tau_{R_0}] = 2 \int_0^{R_0} s'(z) M(z)\,dz
\label{eq:mfpt}
\end{equation}

where $M(z) = \int_0^z m'(y)\,dy$ and $m'(y) = 1/(\sigma^2(y) s'(y))$ is the speed measure density.

We evaluate~\eqref{eq:mfpt} by numerical quadrature in the well-conditioned form

\begin{equation}
\mathbb{E}[\tau_{R_0}] = 2 \int_0^{R_0} g(y)\,dy, \quad g(y) = \frac{1}{\sigma_0^2} \int_1^{R_0/y} s^{-\kappa}\exp(\beta y (s - 1))\,ds
\label{eq:wellcond}
\end{equation}

which avoids the catastrophic cancellation that arises from computing $s(R_0) - s(y)$ directly when $\kappa > 1$ \fref{verification/stochastics/mfpt_bound_explicit.py}.

For the two parameter regimes of interest:

\begin{center}
\small
\begin{tabular}{lrrr}
\toprule
Parameters & Exact MFPT & Bound with $C^*=1$ & Implicit $C^*$ \\
\midrule
$\alpha = 0.5$, $\sigma_0 = 0.8$ & 11.2203 & 19.965 & 0.562 \\
$\alpha = 2.0$, $\sigma_0 = 0.5$ & 443.7374 & $1.415 \times 10^{10}$ & $3.14 \times 10^{-8}$ \\
\bottomrule
\end{tabular}
\end{center}

The values above are the mean first-passage time of the CIR diffusion without jumps, i.e., the $\lambda = 0$ limit of the process defined in Section~\ref{sec:reputation}. The full process with multiplicative downward jumps has a larger MFPT: at $\lambda = 0.8$, the measured value is $19.38$ in the first regime and $710.5$ in the second, and the fraction of paths that reach $R_0$ within the horizon falls from $1.0$ to $0.82$ in the second \fref{verification/stochastics/mfpt_jump_diffusion.py}. The whitewashing bound of this section is therefore a lower bound for the process as defined.

The analytical bound of the original formulation, $C^* \cdot R_0^{1-\kappa} \exp(\beta R_0)$, requires specifying $C^*$. With $C^* = 1$ the bound exceeds the exact MFPT in both regimes: by a factor of 1.78 in the first, and by $3.19 \times 10^7$ (7.5 orders of magnitude) in the second. \textbf{$C^* = 1$ is not a valid lower bound for the CIR first-passage time.} The implicit $C^*$ values (0.562 and $3.14 \times 10^{-8}$) are the largest constants for which the bound holds.

We verified the numerical values by two independent methods. The first is direct Euler simulation of the CIR \fref{verification/stochastics/dt_convergence.py}. The simulated MFPT converges monotonically to the exact value as $dt \to 0$:

\begin{center}
\small
\begin{tabular}{lrr}
\toprule
$dt$ & MFPT ($\alpha = 0.5$, $\sigma_0 = 0.8$) & MFPT ($\alpha = 2.0$, $\sigma_0 = 0.5$) \\
\midrule
0.02   & 12.3115 (+9.7\%)  & 561.55 (+26.6\%) \\
0.005  & 11.7708 (+4.9\%)  & 503.71 (+13.5\%) \\
0.001  & 11.3334 (+1.0\%)  & 461.83 (+4.1\%) \\
0.0005 & 11.2453 (+0.2\%)  & --- \\
\bottomrule
\end{tabular}
\end{center}

The Euler scheme with $\max(R, 0)$ truncation overestimates the MFPT because it delays arrival at the threshold. The bias is monotone in $dt$ and consistent across regimes.

The second method is the censoring test \fref{verification/stochastics/censoring_test.py}, which characterizes the shape of the first-passage time distribution under censoring. In the first regime the process has no censoring within $T_{\max} = 300$: all 20{,}000 simulated paths reach $R_0$. In the second regime only 39.7\% reach $R_0$ at $T_{\max} = 300$. The empirical tail is approximately geometric: the ratio $P_{\text{emp}}/P_{\text{geo}}$ decays from $0.9973$ at $k = 2$ to $0.9858$ at $k = 6$, with mean $0.9927$ and std $0.0038$ across five windows. The deviation from geometricity is small but systematic. Under the geometric assumption the unconditional MFPT would be $T_0/(1-p)$ with $p = P(\tau > T_0) = 0.9214$, giving $635.7$. The exact MFPT is $443.7$. The $43\%$ gap quantifies the effect of the sub-geometric tail: the empirical survival decays faster than geometric, so the extrapolation overestimates the MFPT.

The two methods agree on the qualitative structure: the first-passage time is finite, grows with $R_0$, and is dominated by the ratio $\alpha/\sigma_0^2$. The censoring test is a tail-shape diagnostic, not an independent estimate of the MFPT; its $43\%$ overestimate is the numerical signature of the tail's sub-geometricity. No claim of superexponential growth is made.

\subsection{Stationary Distribution}
Under the Feller condition $2\alpha\bar{R} \geq \sigma_0^2$, the CIR process has a unique stationary distribution:

\begin{equation}
R_\infty \sim \mathrm{Gamma}(k, \theta), \quad k = \frac{2\alpha\bar{R}}{\sigma_0^2}, \quad \theta = \frac{\sigma_0^2}{2\alpha}
\label{eq:gamma}
\end{equation}

with moments

\begin{equation}
\mathbb{E}[R_\infty] = k\theta = \bar{R}, \quad \mathrm{Var}[R_\infty] = k\theta^2 = \frac{\sigma_0^2\bar{R}}{2\alpha}
\label{eq:moments}
\end{equation}

We verify~\eqref{eq:gamma} and~\eqref{eq:moments} by simulating the process from $R(0) = \bar{R}$ over $T = 50$, discarding the first 50\% of samples as burn-in, and comparing empirical moments against the theoretical values \fref{verification/stochastics/stationary_stats.py}. The results:

\begin{center}
\small
\begin{tabular}{lrrcrr}
\toprule
Parameters & $\mathbb{E}[R_\infty]$ theory & empirical & Var theory & Var empirical \\
\midrule
$\alpha = 0.5$, $\sigma_0 = 0.8$ & 1.000000 & 0.998296 (0.17\%) & 0.640000 & 0.637595 (0.38\%) \\
$\alpha = 2.0$, $\sigma_0 = 0.5$ & 1.000000 & 1.000062 (0.006\%) & 0.062500 & 0.063097 (0.96\%) \\
\bottomrule
\end{tabular}
\end{center}

The relative differences are below 1\% in all four quantities. The $Z$-score of the whitewashing target $R_0 = 2.2$ is $1.50\sigma$ in the first regime and $4.80\sigma$ in the second, consistent with empirical estimates ($1.5050\sigma$ and $4.7770\sigma$).

A two-sample Kolmogorov--Smirnov test against $\mathrm{Gamma}(k, \theta)$ does not reject the stationary distribution at 95\%: $D = 0.00354$, $p = 0.557$ (first regime); $D = 0.00479$, $p = 0.201$ (second regime). Because the simulated values are autocorrelated (50{,}000 paths $\times$ 2{,}500 steps), the KS $p$-value is indicative only. The moment comparison is the primary evidence.

\subsection{Integro-Differential Formulation}
For the Fokker--Planck density $p(r, t)$ of the reputation process, the forward equation is:

\begin{equation}
\frac{\partial p}{\partial t} = -\frac{\partial}{\partial r}[\mu(r) p] + \frac{1}{2}\frac{\partial^2}{\partial r^2}[\sigma^2(r) p] + \lambda \int_0^1 \left[\frac{p(r/y, t)}{y} - p(r, t)\right] \nu(dy)
\label{eq:fokker}
\end{equation}

with $\mu(r) = \alpha(\bar{R} - r)$, $\sigma^2(r) = \sigma_0^2 r$, and the jump term written for the multiplicative specification. The domain is $r \in (0, \infty)$, with absorbing boundary at $r = 0$.

The associated value function $C_R(r)$ for reputational collateral satisfies the integro-differential equation:

\begin{equation}
\frac{\partial C_R}{\partial t} + \mu(r)\frac{\partial C_R}{\partial r} + \frac{1}{2}\sigma^2(r)\frac{\partial^2 C_R}{\partial r^2} + \lambda \int_0^1 [C_R(ry) - C_R(r)] \nu(dy) - r C_R = -\pi(r)
\label{eq:integro}
\end{equation}

with boundary condition $C_R(0) = 0$, interpreted in the viscosity sense. Existence and uniqueness follow from standard arguments for integro-PDEs with polynomial growth of $\pi$ \cite{oksendal2007}. The state $r$ is the reputation level on the positive half-line; Section~\ref{sec:payment} evaluates $C_R$ at the Frobenius norm of the reputation tensor, $r = \|T\|_F$, and the reference configuration uses the closed form $C_R(r) = 0.8 \, r / (1 + r)$.

The running reward implied by the closed form follows from \eqref{eq:integro}. With $C_R'(r) = 0.8/(1+r)^2$, $C_R''(r) = -1.6/(1+r)^3$, and, for the reference jump $y = 0.85$ (mean magnitude $0.15$), $C_R(yr) = 0.68\,r/(1+0.85r)$, the stationary operator ($\partial_t C_R = 0$) yields
\begin{equation}
\pi(r) = -\Big[\mu(r)\,\frac{0.8}{(1+r)^{2}} \;-\; 0.8\,\sigma^{2}(r)\,\frac{1}{(1+r)^{3}} \;+\; \lambda\Big(\frac{0.68\,r}{1+0.85r} - \frac{0.8\,r}{1+r}\Big) \;-\; \frac{0.8\,r^{2}}{1+r}\Big],
\label{eq:running-reward}
\end{equation}
with $\mu(r) = \alpha(\bar{R} - r)$ and $\sigma^{2}(r) = \sigma_0^{2} r$. \fref{verification/stochastics/cr_ide_solver.py} evaluates \eqref{eq:running-reward} on a fixed grid: $\pi(r)$ is negative below an admission threshold $r_0 \approx 0.33$ (first regime) and $\approx 0.61$ (second regime), positive and increasing above it, and asymptotically $\pi(r) \sim 0.8\,r$. In the operational range the reward is positive: $\pi(r) > 0$ for every $r \geq 1.0$ in both regimes, so the closed form is consistent with \eqref{eq:integro} wherever the protocol operates.

\subsection{Summary of the Section}
\begin{center}
\resizebox{\textwidth}{!}{%
\begin{tabular}{lll}
\toprule
Claim & Status & Script \\
\midrule
Additive jumps violate positivity     & Verified & \spath{sim_positivity.py} \\
Multiplicative jumps preserve positivity & Verified & \spath{sim_positivity.py} \\
Exact MFPT via numerical integration  & Verified & \spath{mfpt_bound_explicit.py} \\
$C^* = 1$ is not a valid bound        & Verified & \spath{mfpt_bound_explicit.py} \\
Euler bias decreases with $dt$        & Verified & \spath{dt_convergence.py} \\
D4 tail is geometric; S2.2 has no censoring & Verified & \spath{censoring_test.py} \\
Stationary distribution is $\mathrm{Gamma}(k, \theta)$ & Verified & \spath{stationary_stats.py} \\
\bottomrule
\end{tabular}}
\end{center}

The Section's quantitative content is limited to: positivity verification (40.58\% vs 0.0000\%); exact MFPT by numerical integration (11.2203 and 443.7374); implicit $C^*$ per regime (0.562 and $3.14 \times 10^{-8}$); convergent simulated MFPT with $dt$ (11.2453 at $dt = 5 \times 10^{-4}$); stationary moments (relative difference below 1\%).

\section{Tensor Compression}
\label{sec:tensor}

\subsection{Tensor Representation of Reputation}
Reputation is represented as a three-dimensional tensor

\begin{equation}
T \in [0,1]^{C \times D \times H}
\label{eq:tensor}
\end{equation}

where $C$ indexes chains ($C = 20$ in the reference configuration), $D$ indexes contribution domains ($D = 5$), and $H$ indexes time horizons ($H = 5$). Entry $T[c,d,h]$ is the validator's normalized score on chain $c$, in domain $d$, at horizon $h$.

The tensor has $C \cdot D \cdot H = 500$ entries in the reference configuration. This is the object whose norm is compared against the reputational threshold $\Theta_r$ in the zero-knowledge circuit of Section~\ref{sec:zk}.

\subsection{Tucker Decomposition}
The Tucker decomposition of $T$~\cite{tucker1966} is

\begin{equation}
T \approx G \times_0 A^{(0)} \times_1 A^{(1)} \times_2 A^{(2)}
\label{eq:tucker}
\end{equation}

where $G \in \mathbb{R}^{R_1 \times R_2 \times R_3}$ is the core tensor and $A^{(0)} \in \mathbb{R}^{C \times R_1}$, $A^{(1)} \in \mathbb{R}^{D \times R_2}$, $A^{(2)} \in \mathbb{R}^{H \times R_3}$ are the factor matrices. We compute the decomposition via higher-order SVD (HOSVD), which takes the leading singular vectors of each mode unfolding.

For the reference configuration we use $(R_1, R_2, R_3) = (4, 3, 3)$. The storage required for the Tucker representation is

\begin{equation}
R_1 R_2 R_3 + R_1 C + R_2 D + R_3 H = 146 \text{ entries}
\label{eq:tucker-storage}
\end{equation}

versus $C \cdot D \cdot H = 500$ entries for the uncompressed tensor. The storage ratio is $3.4247\times$ (a 70.8\% reduction) \fref{verification/tensor/tucker_analysis.py}.

\subsection{Reconstruction Error}
The HOSVD is exact for tensors of rank at most $(R_1, R_2, R_3)$ and approximate otherwise. We quantify the reconstruction error

\begin{equation}
\varepsilon(T) = \frac{\|T - \hat{T}\|_F}{\|T\|_F}
\label{eq:error}
\end{equation}

for four tensor types at rank $(4,3,3)$:

\begin{center}
\small
\begin{tabular}{lrrr}
\toprule
Tensor type & $\|T\|_F$ & $\|T - \hat{T}\|_F$ & $\varepsilon(T)$ \\
\midrule
Clean low-rank & 73.068 & 0.000 & $4.8 \times 10^{-16}$ \\
Low-rank + noise ($\sigma = 0.05$) & 45.261 & 0.993 & 2.19\% \\
Random (zero mean) & 6.407 & 5.493 & 85.74\% \\
Blocky (community structure) & 6.775 & 3.559 & 52.53\% \\
\bottomrule
\end{tabular}
\end{center}

Three observations follow.

First, the HOSVD implementation is exact: the clean low-rank tensor is reconstructed to machine precision (relative error $4.8 \times 10^{-16}$). This verifies the decomposition itself.

Second, the error for a low-rank tensor corrupted by noise is equal to the noise-to-signal ratio $\|N\|_F / \|T\|_F$. The decomposition captures the underlying low-rank component exactly and fails to capture only the noise, which is by construction not low-rank.

Third, and most importantly: \textbf{HOSVD does not capture block structure.} A tensor whose entries are constant within communities in mode $C$ has mode-1 and mode-2 unfoldings of full rank. The truncation to $R_1 = 4$ rows discards most of the energy in those modes, and the reconstruction error is dominated by truncation rather than by noise. In the reference configuration the blocky tensor has error 52.53\%, comparable to that of a random tensor.

The practical implication is that Tucker compression is efficient for tensors with low-rank structure and inefficient for tensors with community structure. Which regime the reputation tensor occupies is an empirical question that depends on the correlation of validator scores across chains, domains, and horizons. We do not resolve this question in this paper.

\subsection{Error Versus Rank}
For the low-rank noisy tensor, the reconstruction error decreases monotonically with rank:

\begin{center}
\small
\begin{tabular}{lrrr}
\toprule
$(R_1, R_2, R_3)$ & Storage & Ratio & $\varepsilon(T)$ \\
\midrule
$(1,1,1)$ & 31 & $16.13\times$ & 5.02\% \\
$(2,2,2)$ & 68 & $7.35\times$ & 3.24\% \\
$(3,3,3)$ & 117 & $4.27\times$ & 2.32\% \\
$(4,3,3)$ & 146 & $3.42\times$ & 2.19\% \\
$(5,5,5)$ & 275 & $1.82\times$ & 1.52\% \\
\bottomrule
\end{tabular}
\end{center}

The trade-off is between storage reduction and reconstruction fidelity. The choice of rank is a design parameter; the reference configuration $(4,3,3)$ is a specific operating point, not an optimum.

\subsection{Tucker Is Not in the Security Path}
The Tucker compression is a storage and transmission optimization. It is not part of the trust chain of the protocol.

The zero-knowledge proof of Section~\ref{sec:zk} verifies a threshold on the norm of the reputation tensor. The proof operates on the committed tensor, not on its Tucker approximation. A prover who fabricates the factors $G$ and $A^{(i)}$ cannot produce a sound proof of the committed tensor's norm, because the commitment binds the prover to the raw tensor.

Consequently the storage savings of $3.42\times$ do not translate into a reduction of the ZK circuit cost. The circuit must verify the norm of the raw tensor, which is $O(C \cdot D \cdot H)$ constraints under the commitment scheme. Tucker is used only outside the circuit, to reduce the size of the data transmitted from the reputation oracle to the verifier.

This distinction corrects a claim in an earlier formulation of the protocol, in which the storage savings were presented as a reduction of the circuit cost. The two are independent: the storage reduction is real ($3.4247\times$) and the circuit cost reduction is zero. The falsification log records this correction as F-012.

\subsection{Summary of the Section}
\begin{center}
\resizebox{\textwidth}{!}{%
\begin{tabular}{lll}
\toprule
Claim & Status & Script \\
\midrule
HOSVD is exact for low-rank tensors & Verified & \spath{tucker_analysis.py} \\
Storage reduction $500 \to 146$ & Verified & \spath{tucker_analysis.py} \\
Error for noisy low-rank: $\|N\|/\|T\|$ & Verified & \spath{tucker_analysis.py} \\
Error for random: 85.7\% & Verified & \spath{tucker_analysis.py} \\
Error for blocky: 52.5\% (HOSVD fails) & Verified & \spath{tucker_analysis.py} \\
Tucker is not in the security path & Design & Sections~\ref{sec:zk}, \ref{sec:collusion} \\
\bottomrule
\end{tabular}}
\end{center}

The quantitative content of this Section is limited to the storage ratio and the reconstruction errors. The Tucker decomposition is a representation choice, not a security mechanism.

\section{Zero-Knowledge Threshold Proof}
\label{sec:zk}

\subsection{Circuit Specification}
The protocol requires validators to prove that the Frobenius norm of the reputation tensor $T$ exceeds a threshold. The tensor $T$ is public (Section~\ref{sec:datalayer}); the proof exists for succinctness, not for privacy:

\begin{equation}
\|T\|_F \geq \Theta_r
\label{eq:threshold}
\end{equation}

The succinctness is what the proof buys: $T$ lives on the source chain, and the destination chain accepts the constant-size proof instead of the tensor. Section~\ref{subsec:succinctness} quantifies the alternative---transmitting $T$ to the destination chain and computing $\|T\|_F$ there---and shows that it is the more expensive path.

The threshold $\Theta_r$ is the minimum Frobenius norm required for admission. It is a design parameter, not derived from first principles in this paper. The reference configuration uses $\Theta_r = 0.5$ [CONJECTURE].

The natural implementation is a zero-knowledge proof~\cite{gmr1985} over an R1CS circuit that verifies~\eqref{eq:threshold}. The prover commits to $T$ in some form and proves knowledge of a witness satisfying the inequality.

The original specification of the protocol used a Tucker-compressed witness: the prover supplies the core tensor $G$ and factor matrices $A^{(0)}, A^{(1)}, A^{(2)}$, the circuit reconstructs $\hat{T} = G \times_0 A^{(0)} \times_1 A^{(1)} \times_2 A^{(2)}$, and verifies $\|\hat{T}\|_F \geq \Theta_r$. This specification is unsound.

\subsection{Soundness Failure}
The circuit accepts a proof of knowledge of some tensor $\hat{T}$ whose norm exceeds $\Theta_r$. It does not check that $\hat{T}$ equals any committed tensor. A prover with zero reputation can therefore choose $G$ and the factor matrices at will and produce a valid proof for any threshold.

We demonstrate the exploit by constructing a tensor with true norm below the threshold and a fabricated tensor with norm above it \fref{verification/zk/soundness_zk_check.py}:

\begin{center}
\begin{tabular}{lrrl}
\toprule
Tensor & $\|\cdot\|_F$ & Threshold & Accepted \\
\midrule
True reputation & 0.223607 & 0.5 & No \\
Fabricated (prover-chosen) & 1.000000 & 0.5 & Yes \\
\bottomrule
\end{tabular}
\end{center}

The exploit is trivial. Groth16~\cite{groth2016} is sound with respect to the circuit; the circuit verifies a property of a value the prover controls completely. There is no commitment binding the prover to any specific tensor.

\subsection{Repair Options}
Four designs are considered, three of which are sound \fref{verification/zk/soundness_zk_check.py}. We add a fifth column with the on-chain gas cost of verification, measured locally with Foundry \fref{integrations/avalanche/benchmarks/src/GasBenchmarks.sol}. The primitive costs on the EVM are

\begin{center}
\begin{tabular}{lr}
\toprule
Primitive & Gas \\
\midrule
\texttt{ecAdd} (0x06) & 292 \\
\texttt{ecMul} (0x07) & 6{,}142 \\
\texttt{ecPairing} (0x08) & $45{,}000 + 34{,}000 k + 142$ \\
\bottomrule
\end{tabular}
\end{center}

with $k$ the number of pairings. Groth16 verification uses $k = 4$, giving 181{,}142 gas for the pairing step alone. These measurements match EIP-196 and EIP-197 exactly.

\begin{center}
\resizebox{\textwidth}{!}{%
\begin{tabular}{lllrr}
\toprule
Design & Sound & Constraints & Proof size & On-chain gas (direct) \\
\midrule
As specified (Tucker, no commitment) & No & 573 & 0.20 KB & $\sim$206k \\
Option A: Merkle commitment, full opening & Yes & 1{,}080{,}501 & 0.20 KB & $\sim$470M \\
Option A$'$: Poseidon commitment over full tensor & Yes & 240{,}261 & 0.20 KB & $\sim$30M \\
Option B: Pedersen + Bulletproofs & Yes & 0 & 0.72 KB & $\sim$6.5M \\
\bottomrule
\end{tabular}}
\end{center}

The column ``On-chain gas (direct)'' measures the cost of performing the verification work on-chain without a recursive SNARK. For the first row it is the cost of verifying the Groth16 proof of the small circuit ($\sim$206k, constant). For options A, A$'$, and B it is the cost of the on-chain work the SNARK would otherwise do: hashing the tensor (A, A$'$) or performing the Bulletproofs scalar multiplications (B). With recursion, all three options collapse to $\sim$206k on-chain; the column shows why recursion is not optional.

The constraint count for the unsound design decomposes as

\begin{equation}
R_1 R_2 R_3 + R_1 R_2 R_3 + C D H + 1 = 36 + 36 + 500 + 1 = 573
\label{eq:constraints}
\end{equation}

covering the Tucker factor products, the reconstruction of $\hat{T}$, the sum of squares, and the final comparison. This is an order-of-magnitude estimate of the arithmetic complexity; it is not an exact Groth16 constraint count.

\textbf{Option A} commits to $T$ with a Merkle tree of Poseidon hashes and opens all $N = C \cdot D \cdot H$ leaves inside the circuit. The circuit requires 1{,}080{,}501 constraints. Direct on-chain verification of the resulting Groth16 proof costs approximately 206k gas if the circuit fits in one Groth16 proof, which it does not: 1.08M constraints produce a proving key and a proving time that are not practical for on-chain verification [CONJECTURE: the exact key size and proving time depend on the Groth16 implementation and are not measured in this paper]. On-chain direct verification is not viable.

\textbf{Option A$'$} commits to $T$ with a Poseidon tree over the entire tensor inside the circuit. With 240{,}261 constraints the circuit is still too large for practical Groth16 proving on consumer hardware.

\textbf{Option B} replaces the R1CS circuit with a Pedersen commitment to the flattened tensor and a Bulletproofs inner-product argument. The proof size for $n = 500$ elements is $O(\log n) \approx 0.72$ KB. The R1CS constraint count is zero because there is no circuit; soundness is provided by the special-soundness of the Bulletproofs protocol and the binding property of the Pedersen commitment under the discrete logarithm assumption. Direct verification on-chain requires $2N$ scalar multiplications and $2N$ point additions, or approximately 6.5M gas for $N = 500$. This is $31\times$ the cost of Groth16 verification.

None of the three sound designs admits direct on-chain verification at acceptable cost. Option A and Option A$'$ are prohibitively large at the constraint level. Option B is prohibitively expensive at the gas level. The viable path is recursive SNARK composition, treated in Section~\ref{subsec:recursive}.

\subsection{Recursive Verification}
\label{subsec:recursive}
The Bulletproofs verification, which is $O(N)$ in scalar multiplications and costs $\sim$6.5M gas for $N = 500$ if performed on-chain, is instead performed off-chain. A recursive SNARK~\cite{groth2016} then attests that the off-chain verification was correct, and only the recursive proof is posted on-chain.

The recursive proof has constant size and constant verification cost. For a Groth16 recursion over BN254, the on-chain cost is approximately 250k gas, dominated by the \texttt{ecPairing} step at 181{,}142 gas. This is the same order of magnitude as the cost of verifying a single Groth16 proof directly.

We do not measure the proving time of the recursive proof in this paper. The recursive circuit is larger than the inner verification circuit by a factor that depends on the SNARK used for recursion. For the reference configuration, we estimate the recursive proving time in the range of tens of seconds on a modern CPU, but the estimate is not verified.

The choice of recursion framework (Groth16 recursion, Plonk recursion, Halo2 recursion) is not fixed in this paper. All three produce constant-size on-chain proofs at similar gas cost; the difference is the proving time and the trusted setup requirements.

The protocol adopts Option B with recursive SNARK composition. Selecting it does not remove the Tucker representation from the protocol: Tucker remains useful for data transmission outside the circuit (Section~\ref{sec:tensor}), but the circuit operates on the raw committed tensor.

\subsection{Succinctness and the Cost of Ingesting the Tensor}
\label{subsec:succinctness}
The proof of Section~\ref{sec:zk} is sound by commitment binding and constant-size by construction. Because the tensor $T$ is public, a reviewer may ask why a proof is needed at all: a validator could transmit $T$ to the destination chain and compute $\|T\|_F = \sum_i T_i^2$ natively. We compose the cost of this alternative from two terms \fref{verification/zk/succinctness_comparison.py}.

First, the calldata cost of transmitting the tensor (EIP-2028: 4 gas per zero byte, 16 gas per non-zero byte). With $N = C D H = 500$ entries sent as BN254 field elements of 254 bits (one 32-byte word each), the payload is $500 \times 32 = 16{,}000$ bytes, or $256{,}000$ gas at the upper bound. Second, the on-chain sum of squares $\sum_i T_i^2$, measured with Foundry \fref{integrations/avalanche/benchmarks/src/GasBenchmarks.sol}: $428{,}064$ gas at the reference compile, which disables the optimizer exactly as the precompile benchmarks of Section~\ref{subsec:evalgas} do. The composed ingest is therefore $684{,}064$ gas ($455{,}430$ with the optimizer enabled).

This is $3.3\times$ the cost of verifying a Groth16 proof directly ($\sim$206k gas) and $2.7\times$ the cost of the recursive proof ($\sim$250k gas); only the direct on-chain Bulletproofs verification ($\sim$6.5M gas, Section~\ref{subsec:recursive}) is more expensive. Transmitting the tensor and computing its norm on-chain does not dominate the constant-size proof, so the succinctness argument holds on the reference tensor.

\subsection{Estimates and Their Provenance}
The gas measurements for the EVM precompiles are exact: they match EIP-196 and EIP-197 to the gas unit. The Poseidon-2 measurement of 104{,}963 gas \fref{integrations/avalanche/benchmarks/src/GasBenchmarks.sol} is an upper bound from an unoptimized Solidity implementation; production implementations are 20--40k gas. The relevant constraint on Poseidon in this paper is the R1CS count of 240 per hash, which is independent of the on-chain gas cost.

The recursive SNARK cost of $\sim$250k gas is an estimate. The Bulletproofs proof size of 0.72 KB for $n = 500$ is an estimate from the literature. The proving time of the recursive proof is not measured.

The exploit demonstration of Section~\ref{sec:zk} is not an estimate. The values 0.223607 and 1.000000 are computed by the script.

\subsection{Summary of the Section}
\begin{center}
\resizebox{\textwidth}{!}{%
\begin{tabular}{lll}
\toprule
Claim & Status & Script \\
\midrule
Original circuit is unsound & Verified & \spath{soundness_zk_check.py} \\
Fabricated tensor passes threshold & Verified & \spath{soundness_zk_check.py} \\
Constraint count for unsound design & Arithmetic & \spath{soundness_zk_check.py} \\
Merkle opening: $\sim 1.08 \times 10^6$ constraints & Arithmetic & \spath{soundness_zk_check.py} \\
Poseidon tree: $\sim 2.4 \times 10^5$ constraints & Arithmetic & \spath{soundness_zk_check.py} \\
Bulletproofs: 0 R1CS, $\sim$0.72 KB & Estimate & \spath{soundness_zk_check.py} \\
Groth16 $k=4$ pairing: 181{,}142 gas & Verified & \spath{GasBenchmarks.sol} \\
Bulletproofs direct on-chain: $\sim$6.5M gas & Estimated & \spath{GasBenchmarks.sol} \\
Ingesting $T$ on-chain: 684{,}064 gas & Measured & \spath{succinctness_comparison.py} \\
Ingesting $T$ is not cheaper than the proof & Verified & \spath{succinctness_comparison.py} \\
Recursion required for on-chain viability & Design & --- \\
\bottomrule
\end{tabular}}
\end{center}

The Section demonstrates the unsoundness of the original design with a concrete exploit, and recommends a replacement. The exact cost of the replacement depends on implementation parameters not fixed in this paper; the arithmetic is correct given those parameters.

\section{Collusion Detection}
\label{sec:collusion}

\subsection{The Collaboration Graph}
Validators interact through shared governance votes, co-authored pull requests, and simultaneous liquidity events~\cite{kamvar2003,xiong2004}. We model these interactions as a symmetric weighted graph with adjacency matrix

\begin{equation}
W \in \mathbb{R}^{n \times n}
\label{eq:collab}
\end{equation}

where $W_{ij}$ counts the interactions between validators $i$ and $j$. The graph has $n = 500$ nodes in the reference configuration, with mean degree $\bar{d} = 20$.

The hypothesis of interest is that a subset of validators coordinates to inflate each other's reputation without genuine contribution. In graph terms, this appears as a dense subgraph (a clique) embedded in a sparser background.

\subsection{The Normalization Problem}
A natural spectral approach analyzes the symmetric normalized adjacency matrix

\begin{equation}
A_{\mathrm{sym}} = D^{-1/2} W D^{-1/2}
\label{eq:asym}
\end{equation}

where $D = \mathrm{diag}(d_1, \ldots, d_n)$, $d_i = \sum_j W_{ij}$. This matrix has a bulk spectrum that recovers a Wigner-type distribution~\cite{wigner1958,bai2010} for heavy-tailed graphs, which the original formulation of the protocol used for detection.

Each edge is scaled by $1/\sqrt{d_i d_j}$, so the per-edge weight of an intra-clique edge is $1/27.2 \approx 0.037$ and of a background edge is $1/20.0 \approx 0.050$. A per-edge comparison alone would conclude that the clique's edges are weaker after normalization, cancelling the density signal. The relevant comparison is per-pair: intra-clique pairs are $10\times$ more probable than background pairs, so the expected weight is what determines the spectral signature. Quantitatively, for a clique of size $m = 20$ with intra-clique density $10\times$ the background:

\begin{center}
\begin{tabular}{lr}
\toprule
Quantity & Value \\
\midrule
Degree of a clique node & $\approx 27.2$ \\
Degree of a background node & $\approx 20.0$ \\
Normalized intra-clique edge & $1/27.2 \approx 0.037$ \\
Normalized background edge & $1/20.0 \approx 0.050$ \\
\bottomrule
\end{tabular}
\end{center}

The clique's per-edge weight is lower after normalization, but the per-pair weight is not. Under $W$, the expected per-pair weight of an intra-clique pair is $10.0000\times$ that of a background pair. Under $A_{\mathrm{sym}}$, with the null-model background degree cited above ($d_{\mathrm{bg}} = 20$), the ratio is $7.4478$; with the model-consistent background degree under $H_1$ ($d_{\mathrm{bg}} = 27.214$), the ratio is $10.1343$ \fref{verification/rmt/asym_weights.py}. Under either denominator the density signal is attenuated by the normalization, not cancelled. The attenuation under $A_{\mathrm{sym}}$ places the clique signal inside the finite-size bulk edge, while the $10\times$ ratio under $W$ keeps it above; the spectral table below confirms both outcomes.

We verify this numerically \fref{verification/rmt/clique_spectra.py}. For a single clique of size $m = 20$ with intra-clique density $10\times$ the background, the three candidate matrices give the following leading non-trivial eigenvalue, for a single draw of the generator (seed 42):

\begin{center}
\resizebox{\textwidth}{!}{%
\begin{tabular}{llrrrr}
\toprule
Matrix & Statistic & $\lambda$ ($H_0$) & $\lambda$ ($H_1$) & Bulk edge & Gap ($H_1$) \\
\midrule
$W$ (raw) & $s[-2]$ & 8.6896 & 10.6439 & 8.9443 & $+1.6996$ \\
$A_{\mathrm{sym}}$ & $s[-2]$ & 0.4254 & 0.4316 & 0.4472 & $-0.0156$ \\
$B$ & $s[-1]$ & 8.6896 & 10.7011 & 8.9443 & $+1.7569$ \\
\bottomrule
\end{tabular}}
\end{center}

The single-seed values above are one draw from the multi-seed distribution. Across 200 seeds \fref{verification/rmt/clique_spectra_seeds.py}, the $W$ gap has mean $+0.73$ with std $0.39$, the $A_{\mathrm{sym}}$ gap has mean $-0.02$ with std $0.01$, and the $B$ gap has mean $+0.77$ with std $0.41$. The single-seed values sit at the P99, P90 and P99 of their respective distributions. The paired claim $\mathrm{gap}_B > \mathrm{gap}_W$ is structural rather than informative: it follows from Weyl interlacing on $B = W - dd^{\top}/(2m)$. The empirical content is the mean difference $+0.0353$, small against the per-seed standard deviation.

The choice of statistic differs by matrix. For $W$ and $A_{\mathrm{sym}}$, $s[-1]$ is the Perron eigenvalue (trivial, equal to $d_{\max}$ and 1 respectively) and $s[-2]$ is the leading non-trivial direction. For $B$, the identity $B \cdot \mathbf{1} = 0$ places the trivial eigenvalue at the center of the spectrum, and $s[-1]$ is the signal. The change is documented in falsification log entry F-018.

The table shows that $W$ separates the clique from the bulk (gap $+1.6996$ against the asymptotic edge $2\sqrt{\bar{d}} = 8.9443$). $A_{\mathrm{sym}}$ does not (gap $-0.0156$ against its bulk edge $0.4472$). $B$ separates the clique by a wider gap than $W$ ($+1.7569$ vs $+1.6996$) and exceeds both the theoretical bulk edge and the empirically calibrated threshold $q_{99} = 9.0542$ of Section~\ref{subsec:null-distribution}.

The relationship between $B$ and $W$ is explained by a structural identity. For a regular graph with degree $d$, the modularity matrix is $B = A - (d/n) \cdot \mathbf{1}\mathbf{1}^{\top}$, and since $\mathbf{1}$ is orthogonal to every other eigenvector, the spectrum of $B$ equals that of $A$ with the Perron eigenvalue replaced by 0. Consequently $\lambda_{\max}(B) = \lambda_2(W)$ exactly. The Erd\H{o}s--R\'enyi graph used in the null simulation is almost regular (degrees 10--33, mean 19.88, std 4.22), so the identity holds to within $10^{-5}$ in $H_0$. In $H_1$ the clique raises the degrees of its 20 members, breaking the regularity and leaving $B$ slightly ahead of $W$.

This is the content of the falsification log entry F-002.

\subsection{The Modularity Matrix}
The matrix that resolves both problems is the modularity matrix

\begin{equation}
B = W - \frac{dd^{\top}}{2m}
\label{eq:modularity}
\end{equation}

where $d$ is the degree vector and $m = \frac{1}{2}\sum_i d_i$ is the number of edges. The matrix $B$ satisfies

\begin{equation}
B \cdot \mathbf{1} = 0
\label{eq:kernel}
\end{equation}

which removes the degree bias by construction. Its bulk spectrum under the null is centered at zero with edge $2\sqrt{\bar{d}}$ for the degree-corrected stochastic block model.

The modularity matrix is the standard tool for community detection in the stochastic block model~\cite{chung2003,coja2009}. It has the property that the leading eigenvector outside the bulk corresponds to the community structure, not to the degree distribution.

\subsection{Null Distribution and Calibration}
\label{subsec:null-distribution}
We calibrate the null distribution of $\lambda_{\max}(B)$ and $\lambda_{\min}(B)$ by Monte Carlo on Erd\H{o}s--R\'enyi graphs with $n = 500$ and $\bar{d} = 20$ \fref{verification/rmt/calibration_modularity.py}. The calibration uses $N_{H_0} = 2000$ replicates with seed 42.

\begin{center}
\begin{tabular}{lrrr}
\toprule
Statistic & Mean & Std & $q_{99}$ \\
\midrule
$\lambda_{\max}(B)$ (right tail, cliques) & 8.7679 & 0.1115 & 9.0542 \\
$-\lambda_{\min}(B)$ (left tail, bipartite) & 8.8601 & 0.1101 & 9.1371 \\
\bottomrule
\end{tabular}
\end{center}

The empirical thresholds at significance level $\alpha$ are:

\begin{center}
\small
\begin{tabular}{ll}
\toprule
$\alpha$ & Reject $H_0$ if \\
\midrule
0.01 & $\lambda_{\max}(B) > 9.0542$ or $-\lambda_{\min}(B) > 9.1371$ \\
0.05 & $\lambda_{\max}(B) > 8.9663$ or $-\lambda_{\min}(B) > 9.0428$ \\
\bottomrule
\end{tabular}
\end{center}

The dual-tail structure distinguishes clique-based collusion (right tail) from bipartite polarization (left tail). The right-tail threshold is calibrated at $q_{99} = 9.0542$, which is $2.57\sigma$ above the null mean. The left-tail threshold is calibrated at $9.1371$, which is $2.52\sigma$ above its null mean.

The statistic for the right tail is $s[-1]$, the largest eigenvalue of $B$, not $s[-2]$. The change is documented in falsification log entry F-018: for a balanced two-block structure the signal is at $s[-1] = 17.5$ while $s[-2] = 8.6$ sits inside the bulk. The ``skip Perron'' convention that motivates $s[-2]$ in the adjacency matrix does not apply to the modularity matrix, where $B \cdot \mathbf{1} = 0$ places the trivial eigenvalue at the center of the spectrum.

The calibration above is performed on Erd\H{o}s--R\'enyi graphs. The threshold does not transfer to heavy-tailed graphs. The Chung--Lu null with power-law degree distribution places every $H_0$ replicate above the Erd\H{o}s--R\'enyi threshold (empirical false-positive rate $1.0$ at $\alpha = 0.01$ for $\beta \in \{2.1, 2.3, 2.5, 3.0\}$), and per-null recalibration destroys the power against a single planted clique of size $m = 20$ with intra-clique density $10\times$ the background (true-positive rate $0.000$ at $\beta \leq 2.3$, $0.155$ at $\beta = 3.0$) \fref{verification/rmt/calibration_heavy_tail.py}. The bulk edge $2\sqrt{\bar{d}}$ used in the calibration is not the correct edge for Chung--Lu; the entire heavy-tailed null lies above it. The direction of error is anti-conservative: the threshold flags every graph as collusive. This is recorded as falsification log entry F-019. Running the postponed experiment at $n = 5000$ does not repair the calibration (falsification log entry F-020): the Chung--Lu null places every $H_0$ replicate above the recalibrated Erd\H{o}s--R\'enyi threshold ($q_{99} = 9.20$ at $n = 5000$; empirical false-positive rate $1.0$ for $\beta \in \{2.1, 2.3, 2.5, 3.0\}$), and $\beta = 2.1$ remains degree-capped even at that size. A clique-size sweep that holds the clique fraction fixed (falsification log entry F-021) isolates the failure from finite-size and clique-size effects: at the $n = 500$ fraction ($m = 200$, $4\%$ of $n$) the Erd\H{o}s--R\'enyi baseline recovers full power (true-positive rate $1.000$), whereas every Chung--Lu null stays at $0.085$ or below under its own recalibrated threshold. The failure is therefore structural in the heavy-tailed null, not an artifact of $n$ or of $m$: the modularity-spectrum threshold does not transfer to heavy-tailed graphs, and per-null recalibration preserves the false-positive rate only by destroying the power. Detection is not the binding defence in heavy-tailed regimes; admission control there rests on the effective financial floor $\max\{S_{\min}, \Theta - C_R^{\infty}\}$ (Section~\ref{sec:payment}).

\subsection{Detection Power}
\label{subsec:detection-power}
We measure the power of the right-tail test against clique-based collusion by simulating graphs with $k$ planted cliques of size $m$ and intra-clique density $10\times$ the background \fref{verification/rmt/calibration_modularity.py}. The measurement uses $N_{H_1} = 500$ replicates per configuration with Wilson 95\% confidence intervals.

\begin{center}
\small
\begin{tabular}{lrrr}
\toprule
Configuration & $\lambda_{\max}(B)$ mean & TPR @ $\alpha = 0.01$ & 95\% CI \\
\midrule
$k=1$, $m=20$ & 9.7452 & 0.972 & $[0.954, 0.983]$ \\
$k=3$, $m=20$ & 10.2714 & 1.000 & $[0.992, 1.000]$ \\
$k=5$, $m=20$ & 10.5294 & 1.000 & $[0.992, 1.000]$ \\
$k=1$, $m=40$ & 14.2441 & 1.000 & $[0.992, 1.000]$ \\
$k=3$, $m=40$ & 16.2281 & 1.000 & $[0.992, 1.000]$ \\
\bottomrule
\end{tabular}
\end{center}

Detection is uniform across the five configurations. The lowest true positive rate is 0.972, for a single clique of size 20 with intra-clique density $10\times$ the background. All other configurations are detected at 1.000.

The abrupt-transition pattern that appeared with $s[-2]$ was an artifact of the wrong statistic. With $s[-1]$, the smallest clique is detected at 97.2\% and every other configuration at 100\%. There is no intermediate regime in the simulated range.

The single-clique case at $m=20$ remains the tightest: 0.972 with a Wilson 95\% confidence interval of $[0.954, 0.983]$. The interval is above 0.9, so the detection is reliable. Configurations below $m=20$ were not simulated; the boundary at which detection degrades is below the tested range.

\textbf{Comparison with $A_{\mathrm{sym}}$.} The same test run on the normalized adjacency matrix $A_{\mathrm{sym}}$, with the same clique configuration ($k=1$, $m=20$ on $n = 500$, $\bar{d} = 20$), gives a true positive rate of 0.033 at $\alpha = 0.01$ \fref{verification/rmt/archive/calibracion_rmt.py}, measured over $N_{H_1} = 30$ replicates. The modularity matrix $B$ gives 0.972 on the same configuration, measured over $N_{H_1} = 500$ replicates. The ratio is approximately $29.2\times$. This is the numerical basis for the design decision to replace $A_{\mathrm{sym}}$ with $B$, documented as falsification log entry F-002.

The two Wilson 95\% confidence intervals do not overlap: $A_{\mathrm{sym}}$ is $[0.006, 0.167]$ and $B$ is $[0.954, 0.983]$ \fref{verification/rmt/wilson_intervals.py}. The comparison is statistically significant at 95\% despite the small sample size of the $A_{\mathrm{sym}}$ measurement ($N = 30$).

\subsection{The Detection Boundary}
The detection power measured in Section~\ref{subsec:detection-power} is uniform across the simulated range. All five configurations have true positive rates above 0.97. There is no observable transition from partial to full detection within the tested parameter range.

This is a consequence of using $s[-1]$ as the statistic. With the previous statistic $s[-2]$, the pattern appeared to be controlled by the ratio $k \cdot m^2 / n$ and a boundary around $k \cdot m^2 / n \approx 1.5$ was conjectured. That pattern was an artifact: with $s[-1]$ the same configurations all sit above the (much lower) effective boundary.

\textbf{[CONJECTURE]} The detection boundary for the modularity-spectrum test with $s[-1]$ on the degree-corrected stochastic block model lies below the configurations tested here. We do not locate it precisely; the smallest configuration tested (a single clique of size $m=20$) is detected at 0.972, and configurations smaller than $m=20$ were not simulated. A derivation of the boundary from the Kesten--Stigum threshold or from the random matrix literature would be a separate contribution.

Consequently, collusion at the scale tested here is detectable from the collaboration graph alone. Collusion below the (unlocated) boundary would not be. The protocol relies on economic deterrence through slashing (Section~\ref{sec:datalayer}) and reputational collateral destruction (Section~\ref{sec:reputation}) for the sub-boundary regime, which is not characterized in this paper.

\subsection{Limitations}
The calibration is performed on Erd\H{o}s--R\'enyi and Chung--Lu graphs. Real collaboration graphs have additional structure not captured by these models: community hierarchies, temporal correlations between interactions, and directionality (who initiates contact). The empirical quantiles in Section~\ref{subsec:null-distribution} are valid for the calibration distribution; their validity on real graphs is an open question.

The heavy-tailed case is not merely an open question but a negative result. The modularity-spectrum test does not transfer to heavy-tailed collaboration graphs: the failure persists at $n = 5000$, and a clique-size sweep shows it is structural rather than an artifact of $n$ or of the clique size (falsification log entries F-019, F-020, F-021). In those regimes detection is not the binding defence. The protocol's guarantee against a heavy-tailed collusive core rests on the financial floor $\max\{S_{\min}, \Theta - C_R^{\infty}\}$ (Section~\ref{sec:payment}) together with slashing (Section~\ref{sec:datalayer}) and reputational collateral destruction (Section~\ref{sec:reputation}), not on the spectral test.

The left tail ($-\lambda_{\min}$) is measured on disassortative block structures \fref{verification/rmt/bipartite_power.py}. The true positive rate is 0.990--1.000 across block configurations from balanced (250/250) to heavily unbalanced (50/450), with density ratios $\rho_{\text{between}} \in \{2, 5, 10\}$ and $N_{H_1} = 500$ replicates. The left-tail test has power comparable to the right-tail test.

An adversarial validator who knows the test can attempt to avoid detection by keeping the clique size below the (unlocated) boundary, mixing clique members with honest interactions, or avoiding clique structures entirely. The test has power against the naive adversary (uniform cliques and bipartite structures of moderate size); its power against an adaptive adversary below the boundary is not characterized.

\subsection{Summary of the Section}
\begin{center}
\resizebox{\textwidth}{!}{%
\begin{tabular}{lll}
\toprule
Claim & Status & Script \\
\midrule
$A_{\mathrm{sym}}$ masks cliques & Verified & \spath{clique_spectra.py}, \spath{sim_laplacian.py} \\
Modularity $B$ with $s[-1]$ detects cliques & Verified & \spath{calibration_modularity.py} \\
Null $q_{99}$ (right tail, $s[-1]$) $= 9.0542$ & Verified & \spath{calibration_modularity.py} \\
TPR $k=1$, $m=20$ $= 0.972$ & Verified & \spath{calibration_modularity.py} \\
TPR $k=3$, $m=20$ $= 1.000$ & Verified & \spath{calibration_modularity.py} \\
TPR $k=1$, $m=40$ $= 1.000$ & Verified & \spath{calibration_modularity.py} \\
$A_{\mathrm{sym}}$ vs $B$ ratio $= 29.2\times$ & Verified & \spath{wilson_intervals.py} \\
Detection boundary below tested range & Conjecture & --- \\
Left-tail power measured & Verified & \spath{bipartite_power.py} \\
\bottomrule
\end{tabular}}
\end{center}

The Section's quantitative content is limited to the calibration of the null distribution, the power measurements for both tails, and the demonstration that $A_{\mathrm{sym}}$ fails where $B$ succeeds.

\section{Payment Rule}
\label{sec:payment}

\subsection{The Revelation Problem}
A validator is admitted to the validator set if its hybrid collateral exceeds the threshold:

\begin{equation}
\mathrm{CoA}(i) = S_f(i) + C_R(T(i)) \geq \Theta(t)
\label{eq:admission}
\end{equation}

where $S_f$ is the financial stake, $C_R$ is the reputational collateral, and $T$ is the reputation tensor.

The allocation of collateral between its financial and reputational components depends only on $T$. It does not depend on any private information that the validator might hold. This simplifies the mechanism design problem relative to a classic auction, but it also means that the standard tools of auction theory do not apply directly.

We address the two candidate sources of private information below.

\subsection{Verified Reputation}
The reputation tensor $T$ is verifiable on-chain via the zero-knowledge threshold proof of Section~\ref{sec:zk}. A validator whose reputation exceeds the threshold can prove it. A validator whose reputation falls below the threshold cannot produce a valid proof.

Consequently, $T$ is not a private type in the sense of mechanism design. The validator cannot misreport it: the proof binds the validator to a specific tensor, and the verifier checks the binding. This binding property holds under the Option B repair of Section~\ref{subsec:recursive}, which replaces the original circuit of Section~\ref{sec:zk} with a Pedersen commitment and a Bulletproofs argument. The revelation principle does not apply because there is nothing to reveal.

\subsection{Private Cost of Capital}
The validator's cost of capital $\omega \in [\underline{\omega}, \bar{\omega}]$ is private. It represents the opportunity cost of locking financial stake. A validator with low $\omega$ can lock stake cheaply; a validator with high $\omega$ faces a higher opportunity cost.

The allocation rule is

\begin{equation}
S_f^*(T) = \max\{S_{\min}, \Theta(t) - C_R(T)\}
\label{eq:allocation}
\end{equation}

which does not depend on $\omega$. The type affects the validator's utility but not its allocation.

The floor $S_{\min}$ is the technical minimum financial stake, independent of reputation. It is a design parameter, not a value at risk in the financial sense, and it is not arbitrary: it equals the complement of the collateral saturation, $S_{\min} = \Theta - C_R^{\infty}$ with $C_R^{\infty} = \lim_{r \to \infty} C_R(r)$. For the reference value function $C_R(r) = 0.8\,r/(1+r)$ the collateral saturates at $C_R^{\infty} = 0.8$, so $S_{\min} = 1.0 - 0.8 = 0.2$. The floor is the largest minimum stake that stays inactive over the admissible range: the unfloored allocation $\Theta - C_R(T)$ has infimum $\Theta - C_R^{\infty} = 0.2$, which the floor meets but never exceeds \fref{verification/mechanism/financial_floor_sensitivity.py}. A sweep of the saturation $C_R^{\infty}$ over $[0.5, 1.0]$ moves $S_{\min} = \Theta - C_R^{\infty}$ over $[0.0, 0.5]$; the reference value $0.2$ corresponds to $C_R^{\infty} = 0.8$, i.e.\ reputation supplies at most $80\%$ of the admission threshold.

The argument of $C_R$ in~\eqref{eq:allocation} is the Frobenius norm of the tensor: $C_R(T) \equiv C_R(\|T\|_F)$, with $C_R$ the value function of Section~\ref{sec:reputation}.

The utility of a validator of type $\omega$ with reputation $T$ is

\begin{equation}
U(\omega, T) = \pi_{\mathrm{val}} - \omega \cdot S_f^*(T) - F_0
\label{eq:utility}
\end{equation}

where $\pi_{\mathrm{val}}$ is the expected validation revenue per epoch and $F_0$ is the fee charged for participation.

\subsection{The Uniform Posted Price}
Since the allocation does not depend on $\omega$, the only mechanism design decision is the participation fee $F_0$. We choose a uniform posted price: all validators pay the same $F_0$, regardless of type.

The optimal $F_0$ is the largest value that preserves Individual Rationality for all types. The worst-case type is $(\bar{\omega}, T = 0)$, which combines the highest cost of capital with the lowest reputation:

\begin{equation}
U(\bar{\omega}, 0) = \pi_{\mathrm{val}} - \bar{\omega} \cdot \Theta - F_0 \geq 0
\label{eq:ir-worst}
\end{equation}

Solving for $F_0$ at equality:

\begin{equation}
F_0^* = \pi_{\mathrm{val}} - \bar{\omega} \cdot \Theta
\label{eq:f0}
\end{equation}

With the reference parameters ($\pi_{\mathrm{val}} = 0.15$, $\bar{\omega} = 0.10$, $\Theta = 1.0$):

\begin{equation}
F_0^* = 0.15 - 0.10 \cdot 1.0 = 0.0500
\label{eq:f0-value}
\end{equation}

The fee is charged per epoch and is the same for all participating validators.

\subsection{Individual Rationality}
The posted price $F_0^*$ saturates IR at the worst-case type by construction. We verify numerically that IR holds for all types in the support and that the saturation is exact \fref{verification/mechanism/fee_specification.py}.

The Monte Carlo draws $N_{\mathrm{VALIDATORS}} = 1000$ validators with $\omega \sim U[0.01, 0.10]$ and $T$ uniform on $[0, 3]$ (so $C_R(T) \in [0, 0.6]$). The measured utilities are:

\begin{center}
\begin{tabular}{lr}
\toprule
Statistic & Value \\
\midrule
$U_{\min}$ & 0.009747 \\
$U_{\text{mean}}$ & 0.069152 \\
$U_{\max}$ & 0.095891 \\
Fraction with $U \geq 0$ & 1.0000 \\
\bottomrule
\end{tabular}
\end{center}

The analytical saturation $U(\bar{\omega}, 0) = 0$ is exact, with residual below $10^{-9}$. The Monte Carlo $U_{\min}$ is strictly positive (0.009747) because the specific pair $(\omega = 0.10, T = 0)$ is not sampled exactly in a finite draw. The bound holds analytically for the worst-case type and empirically for the sampled population.

\subsection{Sensitivity}
If the fee exceeds $F_0^*$, IR is violated for the worst-case types. We verify this numerically:

\begin{center}
\begin{tabular}{lrrr}
\toprule
$F_0$ & Fraction with $U < 0$ & Gross revenue & Net revenue \\
\midrule
0.025 & 0.0000 & 25.00 & 25.00 \\
0.040 & 0.0000 & 40.00 & 40.00 \\
0.050 & 0.0000 & 50.00 & 50.00 \\
0.060 & 0.0020 & 60.00 & 59.88 \\
\bottomrule
\end{tabular}
\end{center}

At $F_0 = 0.06$ (20\% above $F_0^*$), 0.2\% of the sampled validators violate IR. The gross revenue trade-off is linear: increasing the fee above $F_0^*$ gains revenue at the cost of excluding the worst-case types. Net of the excluded validators, the revenue at $F_0 = 0.06$ is $59.88$ rather than $60.00$.

\subsection{Incentive Compatibility}
Incentive Compatibility is tautological in this mechanism. There is no private type to misreport: $T$ is verified on-chain, and $S_f^*(T)$ does not depend on $\omega$. A validator cannot improve its allocation by misreporting, because the allocation is determined by a public, verifiable signal.

The falsification log records this as F-006. The original formulation of the protocol applied Myerson's optimal auction rule~\cite{myerson1981}, which requires a private one-dimensional type. The current mechanism does not have one. The payment rule is a saturated uniform posted price, not a Myerson auction.

\subsection{Summary of the Section}
\begin{center}
\resizebox{\textwidth}{!}{%
\begin{tabular}{lll}
\toprule
Claim & Status & Script \\
\midrule
$T$ is verified, not private & Design & Section~\ref{sec:zk} \\
$S_f^*(T)$ independent of $\omega$ & Design & --- \\
$S_{\min} = \Theta - C_R^{\infty} = 0.2$ & Analytical & \spath{financial_floor_sensitivity.py} \\
$F_0^* = \pi_{\mathrm{val}} - \bar{\omega} \Theta$ & Analytical & \spath{fee_specification.py} \\
IR holds for all types & Verified & \spath{fee_specification.py} \\
IR saturation exact at worst-case type & Verified & \spath{fee_specification.py} \\
$F_0 = 1.2 F_0^*$ violates IR for 0.2\% & Verified & \spath{fee_specification.py} \\
IC is tautological & Logical & --- \\
\bottomrule
\end{tabular}}
\end{center}

The payment rule is a uniform posted price, not a Myerson auction. The absence of a private type with respect to the allocation rule makes the mechanism simpler than standard auction theory would suggest.

\section{Data Layer}
\label{sec:datalayer}

\subsection{Ingestion of Consensus Events}
The reputation tensor $T$ introduced in Section~\ref{sec:tensor} is not declared by the validator. It is derived from consensus events on the source chains~\cite{saleh2021}. For each chain $c$, domain $d$, and horizon $h$, the entry $T[c,d,h]$ is computed from a formula frozen at genesis:

\begin{equation}
T[c,d,h] = f_c(\text{events in window } h)
\label{eq:ingestion}
\end{equation}

where $f_c$ is a public function that maps the observed consensus events to a normalized score in $[0,1]$. The formula is not changeable by the validator; it is fixed before deployment and applies identically to all validators on chain $c$.

The events themselves are consensus outputs: signed attestations, governance votes, and liquidity events that have been finalized on the source chain. They are not subjective judgments or off-chain declarations.

\subsection{Transport and Freshness}
The commitment to the reputation tensor $T$ is produced by the data layer, not by the validator. The data layer reads the source-chain events, applies the public ingestion formula $f_c$ of Section~\ref{sec:datalayer}, and publishes the commitment optimistically. Soundness rests on the fraud-proof and slashing mechanisms of the data layer, not on the validator's honesty.

For Avalanche L1 source chains, the commitment root is propagated via Avalanche Warp Messaging through the Teleporter contract, with the source chain's quorum as the trust anchor. For non-Avalanche source chains, the transport requires a light client or a zero-knowledge bridge.

The tensor $T$ is frozen per epoch. Two consecutive epochs cannot differ in more than $\gamma \cdot N$ entries, where $N = CDH$ and $\gamma$ is a freshness parameter. The freshness constraint bounds the propagation of undetected fraud across epochs.

\subsection{The Fraud Proof Mechanism}
Any party can submit a fraud proof consisting of two state inclusion proofs against the same source-chain root that yield contradictory values for a single entry $T[c,d,h]$. If the proof is valid, the entry is marked fraudulent and the honest reporter receives a bounty $B$. The validator whose entry was fraudulent is slashed.

We model the audit equilibrium as follows. Let $N_{\mathrm{ATT}}$ be the number of attestations per epoch, $p_{\mathrm{fraud}}$ the fraction that are fraudulent, and $c_{\mathrm{audit}}$ the cost of auditing one attestation. Each verifier decides whether to audit based on the expected revenue~\cite{roughgarden2021}:

\begin{equation}
\mathbb{E}[\text{revenue}] = \frac{p_{\mathrm{fraud}} \cdot N_{\mathrm{ATT}} \cdot B}{n_{\mathrm{auditors}}} - c_{\mathrm{audit}}
\label{eq:revenue}
\end{equation}

where $n_{\mathrm{auditors}}$ is the number of active auditors. A verifier enters if $\mathbb{E}[\text{revenue}] > 0$ and exits otherwise. The equilibrium is

\begin{equation}
n^*(B) = \left\lfloor \frac{B}{B^*} \right\rfloor
\label{eq:equilibrium}
\end{equation}

with the threshold fee

\begin{equation}
B^* = \frac{c_{\mathrm{audit}}}{p_{\mathrm{fraud}}}
\label{eq:bstar}
\end{equation}

For $B < B^*$ no auditor enters; for $B \geq B^*$ at least one auditor is active. The threshold $B^*$ is the fee at which a single auditor is indifferent between auditing and not.

\subsection{Numerical Verification}
We verify~\eqref{eq:equilibrium} and~\eqref{eq:bstar} by Monte Carlo \fref{verification/data-layer/simulate_data_layer_v3.py}. The reference parameters are $N_{\mathrm{ATT}} = 500$, $c_{\mathrm{audit}} = 0.10$, $p_{\mathrm{fraud}} = 0.10$, $V_{\mathrm{fraud}} = 1.0$, $S_{\mathrm{slash}} = 10.0$, $N_{\mathrm{VERIFIERS}} = 20$. The threshold is $B^* = 0.10/0.10 = 1.0$.

\begin{center}
\resizebox{\textwidth}{!}{%
\begin{tabular}{lrrrrr}
\toprule
$B$ & $n^*(B)$ & Detection & Attacker profit & Verifier profit & Integrity \\
\midrule
0.75 & 0 & 0.000 & $+49.91$ & 0.00 & 0.9002 \\
1.00 & 1 & 1.000 & $-502.00$ & $+0.20$ & 1.0000 \\
2.00 & 2 & 1.000 & $-502.00$ & $+0.20$ & 1.0000 \\
5.00 & 5 & 1.000 & $-502.00$ & $+0.20$ & 1.0000 \\
10.00 & 10 & 1.000 & $-502.00$ & $+0.20$ & 1.0000 \\
20.00 & 20 & 1.000 & $-502.00$ & $+0.20$ & 1.0000 \\
\bottomrule
\end{tabular}}
\end{center}

The table reports the interior equilibrium points, where $B$ is either below the threshold or an exact integer multiple of $B^*$. The transition at $B = B^*$ is exact. For $B$ below the threshold no verifier audits; the attacker earns $V_{\mathrm{fraud}} \cdot N_{\mathrm{ATT}} \cdot p_{\mathrm{fraud}}$ per epoch without friction. For $B$ at or above the threshold the auditor coverage reaches 100\% in this model, and the attacker's expected profit is negative because the slashing dominates the fraud gain.

At the interior equilibrium points the analytical verifier profit is zero: the per-auditor revenue $D \cdot B / n^*$ equals the per-auditor cost $N_{\mathrm{ATT}} \cdot c_{\mathrm{audit}} = 50$ exactly when $n^* = B/B^*$. The observed $+0.20$ is the Monte Carlo deviation from zero, equal to $(D_{\text{obs}} - D_{\text{theory}}) \cdot B / n^*$. With $D_{\text{obs}} \approx 50.2$ and $D_{\text{theory}} = 50.0$, the deviation is $0.2 \cdot B^* = 0.20$. Between interior points the verifier profit is strictly positive; the value depends on $B$ and is tabulated in the script output.

\subsection{Slashing and Integrity}
The attacker profit is negative for all $B \geq B^*$ because the slashing parameter $S_{\mathrm{slash}} = 10.0$ exceeds the fraud gain $V_{\mathrm{fraud}} = 1.0$ by an order of magnitude. The design of the slashing parameter is outside the scope of this section; the relevant observation is that once the audit threshold is crossed, the fraud is deterred.

The reputation of a slashed validator is reduced according to the process of Section~\ref{sec:reputation}. The combination of financial slashing and reputation destruction makes fraud strictly unprofitable for any $B$ above the threshold.

\subsection{Freshness as Second Line of Defense}
The freshness parameter $\gamma$ bounds the propagation of undetected fraud across epochs. If a fraudulent entry is present in epoch $t$, it can propagate to epoch $t+1$ only within the $\gamma$ fraction allowed. Over $k$ epochs, the maximum fraction of fraudulent entries in the tensor is bounded by

\begin{equation}
\text{fraud fraction} \leq \gamma \cdot k
\label{eq:freshness}
\end{equation}

in the worst case, where the adversary can insert fresh fraud each epoch. The growth is linear, not exponential. For $\gamma = 0.10$ and $k = 10$ epochs, the maximum fraudulent fraction is 1.0; for $\gamma = 0.05$ the maximum is 0.5 after 10 epochs. The freshness constraint prevents a single undetected fraud from contaminating the entire tensor.

\subsection{Summary of the Section}
\begin{center}
\resizebox{\textwidth}{!}{%
\begin{tabular}{lll}
\toprule
Claim & Status & Script \\
\midrule
$B^* = c_{\mathrm{audit}} / p_{\mathrm{fraud}}$ & Analytical & \spath{simulate_data_layer_v3.py} \\
$n^*(B) = \lfloor B/B^* \rfloor$ equilibrium & Verified & \spath{simulate_data_layer_v3.py} \\
Attacker profit $-502$ with $S_{\mathrm{slash}} = 10$ & Verified & \spath{simulate_data_layer_v3.py} \\
Integrity $= 1.0$ when $n^* \geq 1$ & Verified & \spath{simulate_data_layer_v3.py} \\
Freshness bounds propagation linearly & Analytical & --- \\
\bottomrule
\end{tabular}}
\end{center}

The Section describes the ingestion, transport, and fraud-proof mechanism of the data layer, and verifies the audit equilibrium by Monte Carlo. The freshness parameter provides a second line of defense in case the fraud proof mechanism fails to detect a particular attack.

\section{Bootstrap Dynamics}
\label{sec:bootstrap}

\subsection{The Genesis Problem}
At the genesis block, no validator has reputation. The reputation tensor $T$ is identically zero and the reputational collateral $C_R(0)$ is zero. The protocol therefore begins as a pure proof-of-stake system with full financial stake requirements, and transitions to a hybrid regime as validators accumulate reputation.

The bootstrap question is: does the transition concentrate reputation in the early validators, and if so, how much? A mechanism that concentrates reputation in a small set of genesis validators would recreate the centralization it was designed to avoid.

\subsection{Reputation Models}
The dynamics of reputation during the bootstrap determine the outcome. We compare three models in earlier work and select one for the paper. The models differ in how a validator's activity $a$ translates into reputation growth:

\begin{center}
\small
\begin{tabular}{lll}
\toprule
Model & Gain rule & Stationary value \\
\midrule
Saturating & $k \cdot a \cdot (1-R)$ & $R \to 1$ for all $a$ \\
Decreasing & $k \cdot a \cdot (1-R)^2$ & $R \to 1$ for all $a$ \\
net\_gain & $g \cdot a \cdot (1-R) - \delta R$ & $R^* = g a / (g a + \delta)$ \\
\bottomrule
\end{tabular}
\end{center}

The saturating and decreasing models converge to $R \to 1$ for all validators regardless of activity. After a sufficiently long horizon, the reputation is identical for all validators and carries no information about contribution. This is the content of the earlier comparison \fref{verification/bootstrap/compare_reputation_models.py}.

The net\_gain model has a stationary value that depends on activity:

\begin{equation}
R^* = \frac{g \cdot a}{g \cdot a + \delta}
\label{eq:netgain}
\end{equation}

The bootstrap operates on a bounded reputation score $R \in [0, 1)$, distinct from the unbounded state $r$ of Section~\ref{sec:reputation}; the two are related by $R = r/(1+r)$. During the bootstrap phase (from $R(0) = 0$, before the protocol reaches its stationary regime) the value function $C_R(r)$ of Section~\ref{sec:reputation} is replaced by the bootstrap-specific form $\tilde{C}_R(R) = 0.8 \, R / (1 + R)$. All bootstrap numbers in this section are computed under $\tilde{C}_R$. The value function $C_R(r)$ applies to the post-bootstrap steady state.

The correlation between activity and reputation at stationarity is 0.76 in the reference simulation, versus 0.08--0.11 for the saturating models \fref{verification/bootstrap/compare_reputation_models.py}. We therefore use net\_gain throughout the bootstrap analysis.

The sweep over $\delta$ at fixed $g = 0.01$ \fref{verification/bootstrap/sweep_net_gain.py} gives a maximum of $R_{\mathrm{std}}$ at $\delta/g = 1.5$. We use $g = 0.01$, $\delta = 0.015$.

The quantity $R_{\mathrm{std}}$ is the cross-sectional sample standard deviation of the reputation score $R$ over the active validators at the final epoch of the simulation, with the population convention (numpy default \texttt{ddof = 0}); the reference simulation runs $N = 100$ validators over $T = 500$ epochs \fref{verification/bootstrap/r_std_specification.py}.

\subsection{Bootstrap Simulation}
We simulate the bootstrap from $R(0) = 0$ over 500 epochs \fref{verification/bootstrap/simulate_bootstrap_net_gain.py}. Parameters:

\begin{itemize}[leftmargin=*]
\item $N_{\mathrm{MAX}} = 100$ validators, $N_{\mathrm{INIT}} = 10$ at genesis.
\item Entry rate 0.03 per epoch for new validators.
\item Activity $a \sim \mathrm{LogNormal}(0, 0.5)$ per validator.
\item Transition rate $\tau(t) = \min(1, t/200)$, reaching 1 at epoch 200.
\end{itemize}

We compare three policies:

\begin{itemize}[leftmargin=*]
\item \textbf{pure\_pos}: $S_f = \Theta$ throughout. No reputation adjustment.
\item \textbf{hybrid\_no\_decay}: $S_f(t) = \max\{S_{\min}, \Theta - \tau(t) \cdot C_R(R(t))\}$.
\item \textbf{hybrid\_decay}: same as hybrid\_no\_decay with $R$ decaying at half-life 50 epochs.
\end{itemize}

\subsection{Results}
\begin{center}
\resizebox{\textwidth}{!}{%
\begin{tabular}{lrrrrrr}
\toprule
Policy & $t(\tau{=}1)$ & Gini & Top-10 & $S_f$ med & $R$ med & $R$ max \\
\midrule
pure\_pos & 500 & 0.1283 & 0.1367 & 1.0000 & 0.3936 & 0.5976 \\
hybrid\_no\_decay & 200 & 0.1283 & 0.1367 & 0.7764 & 0.3936 & 0.5976 \\
hybrid\_decay & 200 & 0.1563 & 0.1477 & 0.8393 & 0.2553 & 0.4348 \\
\bottomrule
\end{tabular}}
\end{center}

The $S_f$ med column measures the median financial stake required at the end of the simulation. Under pure\_pos it is $\Theta = 1.0$. Under hybrid\_no\_decay it is 0.7764, a reduction of 22.4\%. Under hybrid\_decay it is 0.8393, a reduction of 16.1\%.

The $R$ med column is the median reputation. Under net\_gain it converges to the stationary value $R^* = g a / (g a + \delta)$. For a validator with median activity $a \approx 1$, $R^* \approx 0.4$, consistent with the measured 0.3936.

The Gini and Top-10 columns are identical between pure\_pos and hybrid\_no\_decay because both policies track the same reputation trajectory $R(t)$. The difference between them is in the financial stake policy $S_f(t)$, not in the reputation. Gini and Top-10 are computed on the distribution of $R$, not on $S_f$.

\subsection{Comparison of Policies}
Two observations follow from the table.

\textbf{Decay does not reduce concentration.} The hybrid\_decay policy has higher Gini (0.1563) than hybrid\_no\_decay (0.1283). This is the opposite of the intuition that decay limits accumulation. The reason is that decay lowers the ceiling ($R_{\max}$ drops from 0.5976 to 0.4348) but also lowers the floor, spreading the distribution. Under hybrid\_no\_decay, the top validators saturate near $R \approx 0.6$ and the distribution compresses; under hybrid\_decay, they fall back to $R \approx 0.43$ and the distribution widens. The falsification log records this as F-015.

\textbf{hybrid\_decay is dominated on both metrics.} It has higher Gini (0.1563 vs 0.1283) and higher median $S_f$ (0.8393 vs 0.7764). The only advantage is lower $R_{\max}$, which could matter if the protocol has a constraint on maximum reputation. Absent such a constraint, hybrid\_no\_decay is the recommended policy.

\subsection{Recommendation}
The protocol uses hybrid\_no\_decay. The financial stake reduction under this policy is 22.4\% at the median validator. This is lower than the 37\% that an earlier analysis reported under the saturating model [see falsification log entry F-016], because the saturating model produces $R \to 1$ for all validators while net\_gain produces $R^* \approx 0.4$.

The reduction is real but modest. It corresponds to a validator with median reputation needing to lock 77.6\% of the stake it would need under pure proof-of-stake. Validators with above-median reputation lock less; validators below median lock more.

\subsection{Summary of the Section}
\begin{center}
\resizebox{\textwidth}{!}{%
\begin{tabular}{lll}
\toprule
Claim & Status & Script \\
\midrule
net\_gain is the selected model & Verified & \spath{compare_reputation_models.py} \\
$\delta/g = 1.5$ maximizes $R_{\mathrm{std}}$ & Verified & \spath{sweep_net_gain.py} \\
$S_f$ median 0.7764 under hybrid\_no\_decay & Verified & \spath{simulate_bootstrap_net_gain.py} \\
22.4\% stake reduction at median & Verified & \spath{simulate_bootstrap_net_gain.py} \\
hybrid\_decay raises Gini & Verified & \spath{simulate_bootstrap_net_gain.py} \\
Hybrid\_no\_decay is recommended & Design & --- \\
\bottomrule
\end{tabular}}
\end{center}

The Section's quantitative content is limited to the bootstrap simulation under the net\_gain model and the comparison of three policies. The saturation of the alternative models is documented as F-016.

The selection of net\_gain is based on its correlation with activity~\cite{cong2021} (0.76 vs 0.08--0.11 for the saturating models), not on satisfying all four criteria of the model comparison. The standard deviation of reputation under net\_gain is 0.0282 in \spath{compare_reputation_models.py} at $\delta/g = 5.0$ with one seed, and 0.0424 in \spath{sweep_net_gain.py} at the sweep optimum $\delta/g = 1.5$, averaged over five seeds; both fall below the 0.05 threshold used in the comparison script. The threshold is a heuristic, not a requirement, and the net\_gain model is chosen for the correlation property, which the saturating models fail by an order of magnitude.

\section{Evaluation}
\label{sec:evaluation}

\subsection{Scope}
This section evaluates three things: the reproducibility of the numerical claims, the measured on-chain gas costs of the cryptographic primitives, and the stability of the conclusions under parameter variation. Two benchmarks are deferred to future work because they require a live network: the proving time of the recursive SNARK, and the latency of cross-chain message propagation.

\subsection{Reproducibility}
Every quantitative claim in Sections~\ref{sec:reputation} through~\ref{sec:bootstrap} is accompanied by a script in \spath{verification/}. The scripts are self-contained: they depend only on NumPy and SciPy, and they declare their seeds and sample sizes at the top of each file. Running \texttt{pip install -r verification/requirements.txt} followed by the commands listed in \spath{verification/README.md} reproduces every number reported in the paper. The README lists the 35 scripts and their expected runtimes.

The project documents its own falsifications in \spath{docs/audit/numerical-verification.md}. As of this writing the log contains 21 entries (F-001 through F-021), each one a claim that was made, tested, and corrected. The log exists because the numerical verification is the primary defense against errors of the kind that would otherwise propagate through the paper unchallenged.

The scripts run in the following time ranges on a 2020-era desktop:

\begin{center}
\small
\begin{tabular}{lr}
\toprule
Script & Time \\
\midrule
\spath{heavy_tail_scaled_clique.py} & $\sim$24 min \\
\spath{heavy_tail_n5000.py} & $\sim$15 min \\
\spath{recalibrate_top.py} & $\sim$10 min \\
\spath{bipartite_power.py} & $\sim$3 min \\
\spath{max_statistic.py} & $\sim$2 min \\
\spath{calibration_modularity.py} & $\sim$400 s \\
\spath{dt_convergence.py} & $\sim$130 s \\
\spath{mfpt_jump_diffusion.py} & $\sim$90 s \\
\spath{calibration_heavy_tail.py} & $\sim$60 s \\
\spath{sim_positivity.py} & $\sim$34 s \\
\spath{clique_spectra_seeds.py} & $\sim$21 s \\
\spath{censoring_test.py} & $\sim$17 s \\
\spath{clique_spectra.py} & $\sim$17 s \\
\spath{stationary_stats.py} & $\sim$12 s \\
\spath{r_std_specification.py} & $\sim$5 s \\
All others & $< 5$ s each \\
\bottomrule
\end{tabular}
\end{center}

Total runtime for the full suite: approximately 67 minutes on a single core. The dominant costs are the heavy-tail recalibrations (\spath{heavy_tail_n5000.py}, $\sim$15 min; \spath{heavy_tail_scaled_clique.py}, $\sim$24 min) and the RMT calibration (\spath{recalibrate_top.py}, $\sim$10 min).

\subsection{On-Chain Gas Costs}
\label{subsec:evalgas}
The cryptographic design of Section~\ref{sec:zk} has an on-chain component. We measure the gas cost of the EVM precompiles that implement the BN254 pairings and scalar multiplications \fref{integrations/avalanche/benchmarks/src/GasBenchmarks.sol}. The measurements match EIP-196 and EIP-197 to the gas unit:

\begin{center}
\small
\begin{tabular}{lrr}
\toprule
Primitive & Measured gas & EIP formula \\
\midrule
\texttt{ecAdd} (0x06) & 292 & $150$ + overhead \\
\texttt{ecMul} (0x07) & 6{,}142 & $6{,}000$ + overhead \\
\texttt{ecPairing} $k=1$ & 79{,}142 & $45{,}000 + 34{,}000 + $ overhead \\
\texttt{ecPairing} $k=2$ & 113{,}142 & $45{,}000 + 68{,}000 + $ overhead \\
\texttt{ecPairing} $k=4$ & 181{,}142 & $45{,}000 + 136{,}000 + $ overhead \\
\bottomrule
\end{tabular}
\end{center}

The overhead is 142 gas in all five measurements, consistent with a single \texttt{staticcall} dispatch. Groth16 verification over BN254 uses $k = 4$ pairings and therefore costs 181{,}142 gas for the pairing step alone, plus approximately 25k gas for the surrounding verification logic. The total is approximately 206k gas, constant in the size of the circuit.

The direct verification of a Bulletproofs proof of Section~\ref{sec:zk} for $n = 500$ elements requires $2n$ scalar multiplications and $2n$ point additions: approximately 6.5M gas. This is $31\times$ the cost of a Groth16 verification. The recursive SNARK of Section~\ref{subsec:recursive} reduces this to the constant cost of the recursive proof, approximately 250k gas.

The measurement uses a local Foundry instance with Solc 0.8.33. The precompile gas schedule is a constant of the EVM protocol, so the measurements apply identically to any EVM-compatible chain including Avalanche C-Chain.

\subsection{Stability of Conclusions}
The conclusions of the paper depend on parameter choices that are fixed at design time. We check that small variations do not change the qualitative outcome.

\textbf{Reputation dynamics (Section~\ref{sec:reputation}).} The domain validity of the multiplicative-jump specification is verified pathwise: 0 of 100{,}000 trajectories leave the positive domain for any parameter choice that satisfies Feller. The mean first-passage time converges to the exact value as $dt \to 0$, with a relative error of 0.22\% at $dt = 5 \times 10^{-4}$. The conclusions do not depend on the discretization.

\textbf{Collusion detection (Section~\ref{sec:collusion}).} The null quantile $q_{99} = 9.0542$ is calibrated with $N_{H_0} = 2000$ replicates. The standard error of the $q_{99}$ estimate is $\sqrt{p(1-p)/N} / f(q_{99})$, evaluated at the empirical $z$-score of $q_{99}$ (2.568), giving 0.017 in $\lambda_{\max}(B)$ units, or 0.186\% of its value. The 95\% confidence interval is approximately $[9.021, 9.087]$. The corresponding theoretical value, using the standard normal density at $z = 2.326$, is 0.103\%. The lowest true positive rate across the five simulated configurations is 0.972 (single clique of size 20), well above the $\alpha = 0.01$ threshold. The qualitative boundary does not depend on the calibration sample size.

\textbf{Data layer (Section~\ref{sec:datalayer}).} The equilibrium $n^*(B) = \lfloor B/B^* \rfloor$ is exact given the parameters. The attacker profit is $-502$ with slashing $S_{\mathrm{slash}} = 10$ and $+49.91$ without detection. The discontinuity at $B = B^*$ is a property of the equilibrium, not of the Monte Carlo.

\textbf{Bootstrap (Section~\ref{sec:bootstrap}).} The stationary reputation under net\_gain is $R^* = ga/(ga + \delta)$. The measured median $R = 0.3936$ matches the analytic $R^* = 0.4$ for unit activity. The choice $\delta/g = 1.5$ maximizes $R_{\mathrm{std}}$ in the sweep; nearby values give similar results.

\subsection{Deferred Benchmarks}
Two benchmarks are not reported in this paper because they require a live network:

\textbf{Recursive SNARK proving time.} The recursive proof of Section~\ref{subsec:recursive} composes a Groth16 verifier inside a larger circuit. The proving time depends on the choice of recursion framework and the size of the inner circuit. This is measurable locally but requires implementing the recursive circuit, which is outside the scope of the current paper.

\textbf{Teleporter latency.} The protocol uses Avalanche Warp Messaging (AWM) through the Teleporter contract to propagate reputation commitments across subnets. The message latency depends on the Avalanche consensus timing and is bounded by the finality time of the source subnet. This is not measurable without a live testnet deployment.

Both benchmarks are declared as future work. The protocols for measuring them are defined in the repository.

\subsection{Summary of the Section}
\begin{center}
\resizebox{\textwidth}{!}{%
\begin{tabular}{lll}
\toprule
Claim & Status & Source \\
\midrule
Verification suite reproduces all claims & Verified & \spath{verification/} \\
21 entries documented & Verified & \spath{docs/audit/} \\
\texttt{ecPairing} $k=4 = 181{,}142$ gas & Measured & \spath{GasBenchmarks.sol} \\
Groth16 verification $\approx 206$k gas & Measured + estimate & \spath{GasBenchmarks.sol} \\
Bulletproofs direct: $\sim$6.5M gas & Estimated & \spath{GasBenchmarks.sol} \\
Recursive SNARK: $\sim$250k gas & Estimated & --- \\
Conclusions stable under parameter variation & Verified & Multiple \\
Recursive proving time & Deferred & --- \\
Teleporter latency & Deferred & --- \\
\bottomrule
\end{tabular}}
\end{center}

The evaluation confirms the reproducibility of the numerical claims and the measured gas costs of the cryptographic primitives. The two deferred benchmarks are declared as future work with the protocols for measuring them.

\section{Limitations}
\label{sec:limitations}

\subsection{Scope}
This section states the limits of the protocol and of the numerical verification. It complements the methodological limitations of Section~\ref{sec:evaluation} (validity of numerical methods, sample sizes, discretization bias) with limits that would persist even if the numerical methods were exact. These are limits of the design itself, not of the simulation.

The distinction matters. If a reviewer asks why the paper does not claim to have solved X, the answer is one of two things: either X was not correctly measured (a limitation of Section~\ref{sec:evaluation}), or X cannot be solved within the design (a limitation of this section).

\subsection{Adversary Model}
The protocol assumes a rational, profit-maximizing adversary. Every security argument in the paper is conditional on this assumption. Specifically:

\textbf{The slashing of Section~\ref{sec:datalayer} deters fraud only if the adversary prefers profit to loss.} A state actor, an ideological opponent, or a competitor with non-monetary utility is not deterred by slashing. The mechanism design does not apply to them.

\textbf{The cost of Sybil reputation is a market cost.} A validator accumulates reputation by producing consensus events on the source chains, and the cost of these events is set by the market for validation on those chains. A Sybil adversary that controls validators on a source chain pays the marginal cost of those validators, not the market cost. The protocol does not model this case; it assumes the adversary does not control the source chains.

\textbf{The data layer bounty $B^*$ of Section~\ref{sec:datalayer} assumes auditors respond to profit.} If no verifier is willing to audit at any bounty, the fraud is not detected. The equilibrium $n^*(B) = \lfloor B/B^* \rfloor$ has a solution only if $B$ is at least $B^*$. Below that threshold, the mechanism provides no detection regardless of the value of the reputation at stake.

These assumptions are standard in the cryptoeconomic literature but they are not universal. The protocol's guarantees do not extend to adversaries that violate them.

\subsection{Collusion Detection Boundary}
The modularity-spectrum test of Section~\ref{sec:collusion} detects cliques and bipartite structures with high power in the simulated range. The smallest configuration tested (a single clique of size $m=20$) is detected at 0.972 with $s[-1]$ as the right-tail statistic.

The detection boundary below which the test loses power was not located. Configurations smaller than $m=20$ were not simulated, and the theoretical Kesten--Stigum threshold for the degree-corrected SBM with the reference parameters ($n = 500$, $\bar{d} = 20$) is below the tested range. The location of the boundary is an open question.

An adaptive adversary who knows the current test may attempt to construct collusion below the (unlocated) boundary. The protocol relies on economic deterrence for this regime rather than on detection.

\subsection{Collusion by Bipartite Polarization}
The left-tail test measures bipartite polarization via $-\lambda_{\min}(B) > 9.1371$. The power of this test was measured \fref{verification/rmt/bipartite_power.py}: true positive rates of 0.990--1.000 across block configurations from balanced (250/250) to heavily unbalanced (50/450), with density ratios $\rho_{\text{between}} \in \{2, 5, 10\}$ and $N_{H_1} = 500$ replicates.

The left-tail test is therefore not a gap. Its power is comparable to the right-tail test. Both tails have been calibrated and measured.

\subsection{Reputation Model Selection}
The bootstrap analysis of Section~\ref{sec:bootstrap} uses the net\_gain model

\begin{equation}
R' = g \cdot a \cdot (1-R) - \delta \cdot R
\label{eq:netgain-limits}
\end{equation}

selected for its correlation between activity and reputation (0.76 vs 0.08--0.11 for the saturating alternatives). The selection does not satisfy all four criteria of the model comparison: the standard deviation of reputation under net\_gain is 0.0424 at the sweep optimum $\delta/g = 1.5$ (five seeds; the $N = 100$ mean is $0.043970$ with 95\% CI $[0.043351, 0.044588]$), below the 0.05 threshold used in the comparison script. The threshold is a heuristic.

If the true reputation dynamics differ from net\_gain --- if the activity of a validator has decreasing returns, if reputation decays with a different functional form, if there are non-linear effects --- the bootstrap results change. The reference stake reduction is 22.4\% at the median validator \fref{verification/bootstrap/simulate_bootstrap_net_gain.py}. The qualitative ranking of policies (hybrid\_no\_decay dominates hybrid\_decay) is robust to small changes in the gain rule; the specific numbers of other rules are not tabulated in this paper.

\subsection{Community Structure and Tucker Compression}
Section~\ref{sec:tensor} shows that the Tucker decomposition does not capture block structure. A tensor whose entries are constant within communities in mode $C$ has mode-1 and mode-2 unfoldings of full rank, and the truncation to $R_1 = 4$ discards most of the energy. The reconstruction error for a blocky tensor is 52.53\%, comparable to that of a random tensor.

If the real reputation tensor has community structure --- if validators cluster by chain, by domain, or by region --- the Tucker compression is inefficient for that structure. Whether the real tensor is low-rank or blocky is not resolved in this paper. It is an empirical question that requires multi-chain reputation data.

\subsection{Calibration on Synthetic Graphs}
The RMT calibration of Section~\ref{sec:collusion} uses Erd\H{o}s--R\'enyi and Chung--Lu graphs. These are standard null models but they do not capture all features of real collaboration graphs: community hierarchies, temporal correlations between interactions, and directionality.

The empirical thresholds ($q_{99} = 9.0542$ for the right tail, $9.1371$ for the left tail) are valid for the calibration distribution. Their validity on real collaboration graphs is an open question. If the real graphs have different spectral properties, the thresholds need recalibration. The recalibration procedure would be the one described in Section~\ref{subsec:null-distribution} (Monte Carlo estimation of the null quantiles), applied to real collaboration graphs instead of synthetic ones.

\subsection{Bootstrap Assumptions}
The bootstrap simulation of Section~\ref{sec:bootstrap} assumes:

\begin{itemize}[leftmargin=*]
\item Activity follows $\mathrm{LogNormal}(0, 0.5)$ per validator.
\item Validators enter at a rate of 0.03 per epoch.
\item The transition rate $\tau(t)$ is linear in time, reaching 1 at epoch 200.
\end{itemize}

These are design choices, not empirical facts. The stationary reputation $R^* = ga/(ga + \delta)$ depends on the activity distribution; if real activity is heavier-tailed or lighter-tailed, $R^*$ changes.

The bootstrap numbers ($S_f$ median 0.7764, $R$ median 0.3936, Gini 0.1283) are specific to the reference parameters. The qualitative ranking of policies (hybrid\_no\_decay dominates hybrid\_decay) is robust; the specific numbers are not.

\subsection{What the Protocol Does Not Claim}
For clarity, the protocol does not claim:

\begin{itemize}[leftmargin=*]
\item Security against a non-rational adversary.
\item Detection of collusion below the (unlocated) Kesten--Stigum boundary for $s[-1]$.
\item Optimality of the Tucker representation for the real reputation tensor.
\item Validity of the RMT thresholds on real collaboration graphs.
\item Independence of the bootstrap results from the activity distribution.
\end{itemize}

Each of these is a direction for future work. The protocols for addressing them are defined in the repository.

\subsection{Summary of the Section}
\begin{center}
\resizebox{\textwidth}{!}{%
\begin{tabular}{lll}
\toprule
Limitation & Category & Section \\
\midrule
Rational adversary assumption & Design & \ref{sec:datalayer}, 10.2 \\
Sybil cost is market cost & Design & 10.2 \\
Bounty threshold $B^*$ & Design & \ref{sec:datalayer} \\
Detection boundary not located & Design & \ref{sec:collusion} \\
Bipartite polarization measured & Verified & \ref{sec:collusion} \\
Net\_gain does not satisfy all criteria & Selection & \ref{sec:bootstrap} \\
Tucker fails for block structure & Method & \ref{sec:tensor} \\
RMT calibrated on synthetic graphs & Method & \ref{sec:collusion} \\
Bootstrap assumes LogNormal activity & Method & \ref{sec:bootstrap} \\
\bottomrule
\end{tabular}}
\end{center}

The methodological limitations (calibration on synthetic data, specific model choices) are separated from the design limitations (non-rational adversary, sub-boundary collusion) because the two require different responses. The methodological limits can be addressed by collecting more data and recalibrating. The design limits require different mechanisms or different claims.

\section{Conclusion}
\label{sec:conclusion}

\subsection{Summary}
The Diamond Framework replaces a fraction of the financial stake required of a proof-of-stake validator with verifiable, non-transferable reputation. The protocol is built from five components, each verified numerically in this paper:

\begin{enumerate}[leftmargin=*]
\item \textbf{Reputation dynamics} (Section~\ref{sec:reputation}). The process is a CIR diffusion with multiplicative downward jumps. It remains in its positive domain by construction. The exact mean first-passage time is computed by numerical integration and confirmed by two independent simulation methods.
\item \textbf{Tensor representation and compression} (Section~\ref{sec:tensor}). The reputation tensor is compressed by Tucker decomposition for storage and transmission. The compression is exact for low-rank tensors and fails for blocky tensors; it is not used in the security path.
\item \textbf{Zero-knowledge threshold proof} (Section~\ref{sec:zk}). The original circuit specification is unsound. The protocol adopts a Pedersen commitment with a Bulletproofs inner-product argument, recursively composed for constant on-chain verification cost. The EVM precompile costs are measured locally and match EIP-196 and EIP-197 exactly.
\item \textbf{Collusion detection} (Section~\ref{sec:collusion}). The modularity matrix $B = W - dd^{\top}/(2m)$ replaces the normalized adjacency matrix $A_{\mathrm{sym}}$. The empirical null quantile at $\alpha = 0.01$ is 9.0542 with $N_{H_0} = 2000$ replicates. The lowest true positive rate across the five simulated configurations is 0.972; the boundary below which detection degrades is below the tested range. The boundary is characterized as a conjecture.
\item \textbf{Payment, data layer, and bootstrap} (Sections~\ref{sec:payment}--\ref{sec:bootstrap}). The payment rule is a saturated uniform posted price, not a Myerson auction. The data layer equilibrium is $n^*(B) = \lfloor B/B^* \rfloor$ with $B^* = c_{\mathrm{audit}}/p_{\mathrm{fraud}}$. The bootstrap under the net\_gain reputation model reduces the median financial stake by 22.4\%.
\end{enumerate}

\subsection{Methodology}
Every quantitative claim in this paper is paired with a script that reproduces it. Claims without scripts are labeled as conjectures. The project maintains a public log of falsifications (\spath{docs/audit/numerical-verification.md}) with 21 entries: each one a claim that was made during development and tested.

The falsification log includes errors of several kinds: a fabricated simulation figure (F-001), a mismeasurement of a gas cost (F-017), a factor of two in an integral (F-011), a rounding rule that shifted a phase transition (F-014), an unverified projection of a multidimensional type (F-004), and an incorrect eigenvalue statistic in the collusion test (F-018). The corrections are documented in the same file as the original claims.

\subsection{Contributions}
The paper contributes:

\begin{enumerate}[leftmargin=*]
\item A cryptoeconomic protocol that uses verifiable reputation as partial collateral, with an explicit adversary model and a stated detection boundary.
\item A numerical verification suite that covers every quantitative claim in the paper, with declared seeds and sample sizes.
\item A falsification log that documents errors of measurement, analysis, and model selection as they were found, with the corrections.
\item Local benchmarks of the EVM precompile costs for the cryptographic primitives used in the protocol.
\end{enumerate}

\subsection{Open Problems}
The limitations section (Section~\ref{sec:limitations}) identifies several open problems that are not resolved in this paper:

\begin{itemize}[leftmargin=*]
\item \textbf{Non-rational adversaries.} The security arguments assume profit-maximizing adversaries. State actors and ideological opponents are outside the model.
\item \textbf{Sub-boundary collusion.} The detection boundary below which the modularity-spectrum test loses power was not located. The protocol relies on economic deterrence for this regime.
\item \textbf{Bipartite polarization.} The left-tail test is calibrated and its power has been measured (0.990--1.000 across tested block configurations).
\item \textbf{Real reputation data.} The Tucker compression efficiency and the RMT calibration thresholds depend on empirical properties of the reputation tensor and the collaboration graph. Neither is available in this paper.
\item \textbf{Recursive proving time.} The recursive SNARK of Section~\ref{subsec:recursive} is designed but not benchmarked.
\item \textbf{Teleporter latency.} Cross-subnet message propagation requires a live testnet deployment and is not measured.
\end{itemize}

\subsection{Closing Remark}
The protocol is in revision. The current document reports the state of the design at the time of writing, with the measurements that have been performed and the claims that remain open. Further development depends on: multi-chain reputation data for calibration, testnet deployment for latency measurement, and external audit for the cryptographic components.

The paper is intended to be readable as a record of what is verified and what is not. The distinction is maintained throughout, and every claim can be traced to a script, a benchmark, or a falsification log entry.

\subsection{Summary of the Section}
\begin{center}
\resizebox{\textwidth}{!}{%
\begin{tabular}{lll}
\toprule
Component & Status & Evidence \\
\midrule
Reputation dynamics & Verified & \spath{stochastics/} \\
Tensor compression & Verified & \spath{tensor/} \\
ZK threshold proof & Verified & \spath{zk/}, \spath{integrations/avalanche/benchmarks/} \\
Collusion detection & Verified & \spath{rmt/} \\
Payment rule & Verified & \spath{mechanism/} \\
Data layer & Verified & \spath{data-layer/} \\
Bootstrap & Verified & \spath{bootstrap/} \\
Reproducibility & Verified & \spath{verification/README.md} \\
Falsification log & 21 entries & \spath{docs/audit/} \\
External audit & Not performed & --- \\
Real-data calibration & Open & --- \\
Testnet benchmarks & Open & --- \\
\bottomrule
\end{tabular}}
\end{center}

The protocol is ready for external audit of the cryptographic components and for testnet deployment of the recursive verification. The remaining open problems are stated in Section~\ref{sec:limitations} and listed in Section~\ref{sec:conclusion}.

\begin{thebibliography}{99}
\bibitem{bai2010} Z.~Bai and J.~W.~Silverstein. \newblock {\em Spectral Analysis of Large Dimensional Random Matrices}. \newblock Springer Series in Statistics. Springer, 2nd edition, 2010.
\bibitem{bauer2012} F.~Bauer and J.~Jost. \newblock Bipartite and neighborhood graphs and the spectrum of the normalized graph Laplacian. \newblock {\em Communications on Pure and Applied Mathematics}, 2012.
\bibitem{chung2003} F.~Chung and L.~Lu. \newblock The average distances in random graphs with given expected degrees. \newblock {\em Proceedings of the National Academy of Sciences}, 99(25):15879--15882, 2002.
\bibitem{coja2009} A.~Coja-Oghlan. \newblock Graph partitioning via adaptive spectral techniques. \newblock {\em Combinatorics, Probability and Computing}, 19(2):227--284, 2010.
\bibitem{cong2021} L.~W.~Cong, Y.~Li, and N.~Wang. \newblock Tokenomics: Dynamic adoption and valuation. \newblock {\em The Review of Financial Studies}, 34(3):1105--1155, 2021.
\bibitem{gmr1985} S.~Goldwasser, S.~Micali, and C.~Rackoff. \newblock The knowledge complexity of interactive proof systems. \newblock {\em SIAM Journal on Computing}, 18(1):186--208, 1989.
\bibitem{groth2016} J.~Groth. \newblock On the size of pairing-based non-interactive arguments. \newblock In {\em Advances in Cryptology -- EUROCRYPT 2016}, pages 305--326. Springer, 2016.
\bibitem{kamvar2003} S.~D.~Kamvar, M.~T.~Schlosser, and H.~Garcia-Molina. \newblock The EigenTrust algorithm for reputation management in P2P networks. \newblock In {\em Proceedings of the 12th International Conference on World Wide Web (WWW)}, pages 640--651, 2003.
\bibitem{myerson1981} R.~B.~Myerson. \newblock Optimal auction design. \newblock {\em Mathematics of Operations Research}, 6(1):58--73, 1981.
\bibitem{oksendal2007} B.~{\O}ksendal and A.~Sulem. \newblock {\em Applied Stochastic Control of Jump Diffusions}. \newblock Universitext. Springer, Berlin, 2007.
\bibitem{roughgarden2021} T.~Roughgarden. \newblock Transaction fee mechanism design. \newblock In {\em Proceedings of the 22nd ACM Conference on Economics and Computation (EC)}, pages 792--792, 2021.
\bibitem{saleh2021} F.~Saleh. \newblock Blockchain without waste: Proof-of-stake. \newblock {\em The Review of Financial Studies}, 34(3):1156--1190, 2021.
\bibitem{tracy1994} C.~A.~Tracy and H.~Widom. \newblock Level-spacing distributions and the Airy kernel. \newblock {\em Communications in Mathematical Physics}, 159(1):151--174, 1994.
\bibitem{tucker1966} L.~R.~Tucker. \newblock Some mathematical notes on three-mode factor analysis. \newblock {\em Psychometrika}, 31(3):279--311, 1966.
\bibitem{wigner1958} E.~P.~Wigner. \newblock On the distribution of the roots of certain symmetric matrices. \newblock {\em Annals of Mathematics}, 67(2):325--327, 1958.
\bibitem{xiong2004} L.~Xiong and L.~Liu. \newblock PeerTrust: Supporting reputation-based trust for peer-to-peer electronic communities. \newblock {\em IEEE Transactions on Knowledge and Data Engineering}, 16(7):843--857, 2004.
\end{thebibliography}

\end{document}
